% 本文件由 LaTeX 排版 Agent 根据有效 translation.json 自动生成。
% author=Apocrypha; work_id=0807; title=尼苛德摩福音

\chapter{尼苛德摩福音}

% structure: p0001 source=h1 target=omitted reason=page-title-already-rendered-by-parent

\Emph{第一部分名为\WorkTitle{比拉多行实}；第二部分名为\WorkTitle{基督下降阴府}。}\par

\Emph{第一部分，第一希腊文版本}\newline{}\Emph{第一部分，第二希腊文版本}\newline{}\Emph{第一部分，拉丁文版本}\newline{}\Emph{第二部分，希腊文版本}\newline{}\Emph{第二部分，第一拉丁文版本}\newline{}\Emph{第二部分，第二拉丁文版本}\par

\SourceRecord{Translated by Alexander Walker.；Ante-Nicene Fathers；Vol. 8.；Edited by Alexander Roberts, James Donaldson, and A. Cleveland Coxe.；Buffalo, NY: Christian Literature Publishing Co.,；1886.；<http://www.newadvent.org/fathers/0807.htm>.}

% structure: p0001 source=h1 target=section reason=page-title

\section{尼苛德摩福音（第一部，希腊文初稿）}

\Emph{关于我等主耶稣基督的记事，作于般雀·比拉多时期}。\par

\Emph{序言}——我，阿纳尼雅，属副总督卫队，谙熟法律，从圣经认识我等主耶稣基督，因信德来到祂面前，并蒙恩领受圣洗圣事。我也查考当时所写关于我等主耶稣基督案件之记事，即犹太人在般雀·比拉多时期所收藏者，发现这些用希伯来文写成的记事，并蒙天主恩宠，将其译为希腊文，为所有呼求我等师尊耶稣基督之名者提供信息。时值我主弗拉维·狄奥多西在位第十七年，弗拉维·瓦伦提尼安在位第六年，第九小纪。\par

因此，凡阅读并将之转抄于其他书籍者，请记念我，为我祈祷，愿天主怜悯我，赦免我所犯得罪祂的罪过。\par

愿平安归于阅读者、听闻者及其家人。阿们。\par

在提贝里乌斯·凯撒任罗马皇帝第十五年，黑落德为加里肋亚王第十九年，四月朔日前第八日，即三月二十五日，鲁福斯与鲁贝利奥执政之年，第二百零二届奥运会第四年，若瑟·盖法任犹太人大司祭之时。\par

以下为尼苛德摩在我等主耶稣基督、救主天主受十字架苦难之后，用希伯来文所写，并为后世所留的记述：——\par

% structure: p0008 source=h2 target=subsection reason=relative-heading-rank; base=h2

\subsection{第一章}

大司祭和经师们召集公议会，亚纳斯、盖法、塞梅斯、达塔厄斯、加玛里耳、犹达斯、肋未、纳斐塔里、亚历山大、雅依洛，以及其余犹太人，来到比拉多面前，在许多事上控告耶稣，说：我们知道这人是若瑟木匠之子，生于玛利亚；祂却自称是天主子，又是君王；此外，祂亵渎安息日，并想废除我们祖传的法律。比拉多说：祂做了什么事，表明祂想废除法律？犹太人说：我们有法律，不可在安息日医治任何人；但此人在安息日治好了瘸子、驼背、枯干手、瞎子、瘫痪者、哑巴和附魔者，用的是邪术。比拉多对他们说：什么邪术？他们对他说：祂是个术士，靠着魔王贝尔则步驱魔，万物都服从祂。比拉多对他们说：这不是凭着不洁之神驱魔，而是凭着医神\OriginalTerm{阿斯克勒庇俄斯}{Aesculapius}。\par

犹太人对比拉多说：我们恳请大人，让祂站在你的审判席前受审。比拉多召了他们来，说：告诉我，我这个总督怎能审判一个君王？他们对他说：我们不是说祂是君王，而是祂自称是君王。比拉多召来跑差，对他说：把耶稣恭敬地带进来。跑差出去，认出祂，就朝拜了祂，并把自己的外衣拿在手里，铺在地上，对祂说：我主，请从这上面走进来，因为总督召你。犹太人看见跑差所做的事，就对比拉多大喊说：你为什么命令用跑差带祂进来，而不是用传令官？因为那跑差一见祂就朝拜了祂，并把他的上衣铺在地上，让祂像君王一样走过去。\par

比拉多召来跑差，对他说：你为什么这样做，把你的外衣铺在地上，让耶稣从上面走？跑差对他说：我主总督，当你派我去耶路撒冷见亚历山大时，我看见祂骑在一头驴上，希伯来的孩子们手里拿着树枝，呼喊；另有人把自己的衣服铺在祂下面，说：\BibleQuote{贺三纳于至高之天！奉主名而来的，当受赞美。}\BibleRef{玛 21:8}\par

犹太人大喊，对跑差说：希伯来的孩子们用希伯来语呼喊；你从哪里知道希腊语？跑差对他们说：我问了一个犹太人，说：他们用希伯来语喊什么？他就为我翻译了。比拉多对他们说：他们用希伯来语喊什么？犹太人对他说：\OriginalTerm{和散那 梅姆布罗梅 巴鲁恰玛 阿多奈}{Hosanna membrome Baruchamma Adonaï}。比拉多对他们说：这\Quote{和散那}等等，如何解释？犹太人对他说：\BibleQuote{现在拯救，在至高之天；奉主名而来的，当受赞美。}比拉多对他们说：如果你们承认孩子们所说的话，跑差又做错了什么？他们沉默不语。总督对跑差说：出去，随你用什么方式带祂进来。跑差出去，像先前一样做了，并对耶稣说：我主，请进来，总督召你。\par

耶稣进来，执旗者举着旗帜，旗杆顶端自动弯下，朝拜了耶稣。犹太人看见旗帜的姿势，如何弯下朝拜耶稣，就严厉斥责执旗者。比拉多对犹太人说：你们不奇怪旗杆顶端自动弯下朝拜耶稣吗？犹太人对比拉多说：我们看见执旗者弯下旗杆朝拜了祂。总督召来执旗者，对他们说：你们为什么这样做？他们对比拉多说：我们是希腊人，是神庙的仆役，我们怎会朝拜祂？的确，我们举着旗杆时，顶端自动弯下，朝拜了祂。\par

比拉多对会堂首领和民间长老说：\Quote{你们为自己挑选强壮的男子，让他们举着旗帜，我们看看旗帜是否会随他们弯下。}犹太长老选了十二个强壮有力的男子，让他们六人一组举着旗帜；他们被安置在总督审判台前。比拉多对差役说：\Quote{把他带出总督府，然后随你心意再带他进来。}耶稣和差役走出总督府。比拉多召来先前举旗帜的人，对他们说：\Quote{我以凯撒的性命起誓，如果耶稣进来时旗帜不弯下，我就砍你们的头。}总督命耶稣第二次进来。差役像先前一样，再三恳求耶稣踩在他的外衣上。耶稣踩在上面，走了进去。他进去时，旗帜再次弯下，朝拜了耶稣。\par

% structure: p0015 source=h2 target=subsection reason=relative-heading-rank; base=h2

\subsection{第二章}

比拉多见此，害怕起来，想离开审判台；但他正在盘算离开时，他的妻子派人来对他说：\Quote{不要插手这件正义的事，因为今夜我因他受了许多苦。}\BibleRef{玛 27:19}比拉多召来犹太人，对他们说：\Quote{你们知道我的妻子是敬畏天主的人，并且与你们一样信奉犹太教。}他们对他说：\Quote{是的，我们知道。}比拉多对他们说：\Quote{看，我的妻子派人来对我说：\InnerQuote{不要插手这件正义的事，因为今夜我因他受了许多苦。}}犹太人回答比拉多说：\Quote{我们不是告诉过你，他是个术士吗？看，他给你妻子送了一个梦。}\par

比拉多召来耶稣，对他说：\Quote{这些人作证反对你，你有什么可说的？}耶稣说：\Quote{除非他们有能力，否则他们什么也不会说；因为每个人都有权用自己的口说好说歹。他们自己会看着办。}\par

犹太长老回答耶稣说：\Quote{我们要看什么？首先，你是在奸淫中生的；其次，你在白冷的出生导致了婴孩的被杀；第三，你的父亲若瑟和你的母亲玛利亚逃往埃及，因为他们对人民没有信心。}\par

旁观者中有些虔诚的犹太人说：\Quote{我们否认他是奸淫所生；因为我们知道若瑟聘了玛利亚，他不是奸淫所生的。}比拉多对说他是奸淫所生的犹太人说：\Quote{你们的话不真实，因为他们已订婚，正如这些同胞所说的。}亚纳斯和盖法对比拉多说：\Quote{我们全都呼喊他是奸淫所生，却不被相信；这些人是归依者，是他的门徒。}比拉多叫来亚纳斯和盖法，问他们：\Quote{什么是归依者？}他们对他说：\Quote{他们是希腊人的子孙，如今成了犹太人。}而说耶稣不是奸淫所生的那些人，即拉匝禄、阿斯泰里、安东尼、雅各伯、阿姆内斯、泽拉斯、撒慕尔、依撒格、丕乃哈斯、克黎斯普、阿格黎帕和犹达，说：\Quote{我们不是归依者，而是犹太人的子孙，讲的是真理；因为我们当时在场，见证了若瑟和玛利亚的订婚。}\par

比拉多叫来这十二个说耶稣不是奸淫所生的人，对他们说：\Quote{我以凯撒的性命起誓，你们要告诉我，你们所说他没有奸淫所生是否属实。}他们对比拉多说：\Quote{我们有法律禁止起誓，因为那是罪；但他们若以凯撒的性命起誓，说事实并非如我们所说，我们就要被处死。}比拉多对亚纳斯和盖法说：\Quote{你们对此没有什么可回答的吗？}亚纳斯和盖法对比拉多说：\Quote{这十二个人说他没有奸淫所生，他们被相信；但我们全体都呼喊他是奸淫所生，是术士，还说自己是天主子、是君王，我们却不被相信。}\par

比拉多命令所有群众出去，只留下那十二个说耶稣不是奸淫所生的人，并命耶稣与他们分开。比拉多问他们：\Quote{他们为什么想要处死他？}他们对他说：\Quote{他们发怒，因为他在安息日治病。}比拉多说：\Quote{因为一件善事，他们就要处死他吗？}他们对他说：\Quote{是的。}\par

% structure: p0022 source=h2 target=subsection reason=relative-heading-rank; base=h2

\subsection{第三章}

比拉多满腔怒火，走出总督府，对他们说：\Quote{我以太阳为证，我在这个人身上找不到任何过错。}犹太人回答总督说：\Quote{如果这个人不是作恶的，我们就不会把他交给你。}比拉多说：\Quote{你们带他去，按你们的法律审判他。}犹太人对比拉多说：\Quote{我们无权处死任何人。}比拉多说：\Quote{难道是天主说你们不能处死人，而我却能？}\par

比拉多又走进总督府，私下对耶稣说：\Quote{你是犹太人的君王吗？}耶稣回答比拉多说：\Quote{这话是你自己说的，还是别人对你说的关于我的话？}比拉多回答耶稣说：\Quote{难道我是犹太人吗？你的民族和司祭长把你交给了我。你做了什么？}耶稣回答说：\Quote{我的国不属于这世界；如果我的国属于这世界，我的臣仆就会为我战斗，使我不致被交给犹太人；但如今我的国不在这里。}比拉多对他说：\Quote{那么，你是君王了？}耶稣回答他说：\Quote{你说我是君王。我为此而生，为此而来，为叫每一个属于真理的人听我的声音。}比拉多对他说：\Quote{什么是真理？}耶稣对他说：\Quote{真理来自天上。}比拉多说：\Quote{真理不在地上吗？}耶稣对比拉多说：\Quote{你看见那些讲真理的人是如何被地上有权势者审判的。}\par

% structure: p0025 source=h2 target=subsection reason=relative-heading-rank; base=h2

\subsection{第四章}

比拉多把耶稣留在总督府内，自己出去对犹太人说：\Quote{我在他身上找不到任何过错。}犹太人对他说：\Quote{他说：\InnerQuote{我能拆毁这座圣殿，并在三天内重建。}}比拉多说：\Quote{什么圣殿？}犹太人说：\Quote{就是撒罗满用了四十六年建成的圣殿，\BibleRef{若 2:20}而这个人说要拆毁它，并在三天内重建。}比拉多对他们说：\Quote{对这义人的血，我是无辜的。你们自己看着办吧。}犹太人说：\Quote{他的血归在我们和我们子孙身上。}\par

比拉多召集了长老、司铎和肋未人，私下对他们说：\Quote{不要这样做，因为你们对他的指控没有一项是该死的；你们的指控不过是关于治病和安息日的亵渎。}长老、司铎和肋未人说：\Quote{如果有人诽谤凯撒，他该不该死？}比拉多说：\Quote{他该死。}犹太人对比拉多说：\Quote{如果有人诽谤凯撒，他就该死；但这个人诽谤了天主。}\par

总督命令犹太人离开总督府；然后召来耶稣，对他说：\Quote{我该对你怎么办？}耶稣对应比拉多说：\Quote{正如已经赐给你的那样。}比拉多说：\Quote{怎么赐给的？}耶稣说：\Quote{梅瑟和先知们早已预言了我的死亡和复活。}犹太人注意到这一点，听见了，对比拉多说：\Quote{你还要听这亵渎的话到什么时候？}比拉多对犹太人说：\Quote{如果这些话是亵渎的，你们就把他当作亵渎者带走，带到你们的会堂，按你们的法律审判他。}犹太人对比拉多说：\Quote{我们的法律规定，凡损害同人的，该受四十减一的鞭打；但亵渎天主的，该用石头砸死。}\par

比拉多对他们说：\Quote{你们把他带走，随你们怎么处治他。}犹太人对比拉多说：\Quote{我们愿意他钉十字架。}比拉多说：\Quote{他不该钉十字架。}\par

总督环视站在旁边的犹太人群，看见许多犹太人哭泣，就说：\Quote{全体民众都不愿意他死。}犹太人的长老说：\Quote{正因如此，我们全体民众都来了，要他死。}比拉多对犹太人说：\Quote{他为什么该死？}犹太人说：\Quote{因为他自称是天主子、君王。}\par

% structure: p0031 source=h2 target=subsection reason=relative-heading-rank; base=h2

\subsection{第五章}

有一个犹太人尼苛德摩站在总督面前，说：\Quote{求大人，让我说几句话。}比拉多说：\Quote{说吧。}尼苛德摩说：\Quote{我对长老、司铎、肋未人和会堂中所有的犹太群众说过：你们想对这个人做什么？这个人行了许多奇迹和奇事，是没有人行过、也没有人能行的。放了他吧，不要对他存任何恶意。如果他行的奇迹来自天主，它们必会存立；如果出于人，它们必要消散。\BibleRef{宗 5:38} 确实，梅瑟奉天主差遣进入埃及，行了天主吩咐他在埃及王法郎面前行的许多奇迹。那里有法郎的仆人雅乃斯和杨布勒，他们也行了梅瑟所行的一些奇迹；埃及人把他们当作神——就是这雅乃斯和杨布勒。\BibleRef{弟后 3:8} 但是，由于他们所行的奇迹不是出于天主，他们和信从他们的人都灭亡了。现在，释放这个人吧，因为他是不该死的。}\par

犹太人对尼苛德摩说：\Quote{你成了他的门徒，所以替他辩护。}尼苛德摩对他们说：\Quote{也许总督也成了他的门徒，因为他替他辩护。难道皇帝没有委派他担任这显职吗？}犹太人极其愤怒，向尼苛德摩咬牙切齿。比拉多对他们说：\Quote{你们听见了真理，为什么向他咬牙切齿？}犹太人对尼苛德摩说：\Quote{愿你得到他的真理和他的份。}尼苛德摩说：\Quote{阿们，阿们，愿我如你们所说得到它。}\par

% structure: p0034 source=h2 target=subsection reason=relative-heading-rank; base=h2

\subsection{第六章}

有一个犹太人走上前来，请求总督允许他说一句话。总督说：\Quote{你若有什么话，就说吧。}那犹太人说：\Quote{我躺在床上有三十八年，受极大的痛苦。耶稣来了，许多附魔的人和患各种疾病的人都被他治好了。几个青年人可怜我，把我连床带人抬到耶稣那里。耶稣看见我，就怜悯我，对我说：\InnerQuote{拿起你的床走吧。}我就拿起我的床走了。}犹太人对比拉多说：\Quote{问他是在哪一天被治好的。}那被治好的人说：\Quote{是在安息日。}\BibleRef{若 5:5} 犹太人说：\Quote{这不正是我们所说的：他在安息日治病、驱魔吗？}\par

另一个犹太人上前来说：\Quote{我生来是瞎子，能听见声音，却看不见面容。耶稣路过时，我大声呼喊：\InnerQuote{可怜我吧，达味之子！}他就怜悯我，把手放在我的眼睛上，我立即重见光明。}另一个犹太人上前来说：\Quote{我原是伛偻的，他用一句话治直了我。}另一个说：\Quote{我原是癞病人，他用一句话治愈了我。}\par

% structure: p0037 source=h2 target=subsection reason=relative-heading-rank; base=h2

\subsection{第七章}

有一个妇人在远处呼喊说：\Quote{我患血漏，摸了祂的衣边，我十二年的血漏立刻就止住了。}\BibleRef{玛 9:20} 犹太人说：\Quote{我们有法律，不接受女人的见证。}\par

% structure: p0039 source=h2 target=subsection reason=relative-heading-rank; base=h2

\subsection{第八章}

另有一大群男男女女呼喊说：\Quote{这人是先知，连魔鬼都服从他。}比拉多对说魔鬼服从祂的人说：\Quote{那么，你们的老师为什么不也服从祂呢？}他们对比拉多说：\Quote{我们不知道。}又有人说：\Quote{祂使死了四天的拉匝禄从坟墓中复活了。}\BibleRef{若 11:1} 总督战栗起来，对全体犹太群众说：\Quote{你们为什么想要流无辜者的血？}\par

% structure: p0041 source=h2 target=subsection reason=relative-heading-rank; base=h2

\subsection{第九章}

比拉多召来尼苛德摩和那十二个说他不是从淫乱而生的人，对他们说：\Quote{我该怎么办，因为民众骚动了？}他们对他说：\Quote{我们不知道，让他们自己看着办。}比拉多又召集全体犹太群众，说：\Quote{你们知道，在无酵节，按惯例我要释放一个囚犯给你们。狱中有一个已判罪的囚犯，是个杀人犯，名叫巴辣巴；而站在你们面前的这个人——耶稣，我查不出祂有什么罪。你们要我释放哪一个给你们？}他们喊说：\Quote{巴辣巴。}比拉多说：\Quote{那么，对那称为基督的耶稣，我们该怎么办？}犹太人说：\Quote{钉祂在十字架上！}另有人说：\Quote{你若释放这人，就不是凯撒的朋友，因为祂自称是天主子和君王。那么，你要让这个人作王，而不要凯撒吗？}\par

\Person{比拉多}怒气冲冲地对犹太人说：\Quote{你们的民族总是叛逆的，你们总是说反对你们恩人的话。}犹太人说：\Quote{什么恩人？}他对他们说：\Quote{你们的天主带领你们出离\Place{埃及}地，脱离苦役，平安地经过海中如同经过干地，在旷野用玛纳喂养你们，赐给你们鹌鹑，用磐石中的水解了你们的渴，并赐给你们法律；在这一切事上，你们激怒了你们的天主，并寻求铸造一只牛犊。你们使你们的天主恼怒，他想要杀死你们。\Person{梅瑟}为你们祈祷，你们才没有被处死。现在你们却指控我憎恨皇帝。}\par

他从审判台上站起来，想要走出去。犹太人大声喊道：\Quote{我们知道\Person{凯撒}是王，而不是\Person{耶稣}。因为贤士们确实带了礼物来给他，如同给一位王。当\Person{黑落德}从贤士那里听说有一位王诞生了，就想要杀害他；他的父亲\Person{若瑟}知道了这事，就带着他和他的母亲，逃往\Place{埃及}去了。\Person{黑落德}听说此事，就杀害了在\Place{白冷}出生的希伯来人的婴孩。}\par

\Person{比拉多}听了这些话，就害怕起来；他命令群众静默，因为他们正在喊叫，对他们说：\Quote{这就是\Person{黑落德}所寻找的人吗？}犹太人说：\Quote{是的，就是他。}然后，\Person{比拉多}拿了水，当着太阳的面洗手，说：\Quote{对这义人的血，我是无罪的；你们自己看着办吧。}犹太人又喊道：\Quote{他的血归在我们和我们子孙的身上！}\par

于是\Person{比拉多}命令把他坐着的审判台的幔子拉上，对\Person{耶稣}说：\Quote{你的民族控告你是王。为此，我判决你：先受鞭打，按照可敬的君王的法令，然后被钉在十字架上，在你被捕的那个园子里。并且让\Person{狄斯马斯}和\Person{格斯塔斯}这两个凶犯与你一同被钉十字架。}\par

% structure: p0047 source=h2 target=subsection reason=relative-heading-rank; base=h2

\subsection{第十章。}

\Person{耶稣}从总督府出来，那两个凶犯与他同行。他们到了那地方，就脱去他的衣服，给他围上一条毛巾，把一顶荆棘冠冕戴在他头上。他们把他钉在十字架上；同时，也把两个凶犯与他一同挂起来。\Person{耶稣}说：\Quote{圣父，宽恕他们吧，因为他们不知道自己在做什么。}士兵分了他的衣服；民众站着观望他。司祭长和他们的首领们一同讥笑他，说：\Quote{他救了别人，让他救自己吧。如果他是天主子，让他从十字架上下来吧。}士兵也戏弄他，上前来，用醋和苦胆混合的酒给他喝，说：\Quote{你是犹太人的王，救自己吧。}\par

\Person{比拉多}宣判之后，下令将所控告他的罪状用希腊文、拉丁文和希伯来文写成标题，按照犹太人所说的：\Quote{他是犹太人的王。}\par

其中一个被挂起来的凶犯对他说：\Quote{如果你是基督，救你自己和我们吧。}\Person{狄斯马斯}回答说，责备他说：\Quote{你既受同样的刑罚，还不怕天主吗？我们实在是公道的，因为我们所受的与我们的行为相称；但这个人没有做过一件坏事。}他对\Person{耶稣}说：\Quote{主啊，当你的王国来临时，请纪念我。}\Person{耶稣}对他说：\Quote{阿们，阿们，我告诉你：今天你就要同我在乐园里。}\par

% structure: p0051 source=h2 target=subsection reason=relative-heading-rank; base=h2

\subsection{第十一章。}

大约第六时辰，黑暗笼罩大地直到第九时辰，日头变暗；圣殿的幔子从上到下裂为两半。\Person{耶稣}大声呼喊说：\Quote{圣父，\OriginalTerm{巴达赫 厄弗基德 鲁埃尔}{Baddach ephkid ruel}，意即：\InnerQuote{我将我的灵魂交在你手中。}}说了这话，就断了气。百夫长看见所发生的事，就光荣了天主，说：\Quote{这真是个义人。}所有在场的群众，看见所发生的事，都捶着胸膛回去了。\par

百夫长把所发生的事报告给总督。总督和他的妻子听了，极其忧伤，那天没有吃喝。\Person{比拉多}派人叫犹太人来，对他们说：\Quote{你们看见所发生的事了吗？}他们说：\Quote{是日蚀，照常发生了。}\par

他的相识们远远站着，那些从\Place{加里肋亚}同他来的妇女也看见了这些事。一个名叫\Person{若瑟}的人，来自\Place{阿黎玛特雅}城的议员，他也是期待天主之国的，去见\Person{比拉多}，求\Person{耶稣}的遗体。他取下来，用干净的殓布裹好，安放在一座从岩石中凿出的坟墓里，那里从未葬过人。\par

% structure: p0055 source=h2 target=subsection reason=relative-heading-rank; base=h2

\subsection{第十二章。}

\Person{犹太人}听说\Person{若瑟}求得了\Person{耶稣}的遗体，就寻找他和那十二个曾说耶稣不是由淫乱而生的人，以及\Person{尼苛德摩}，还有许多在\Person{比拉多}面前挺身而出、宣告耶稣善行的人。在这些躲藏起来的人中，只有\Person{尼苛德摩}被他们看见，因为他是犹太人的首领。\Person{尼苛德摩}对他们说：\Quote{你们怎么进了会堂？}犹太人对他说：\Quote{你怎么进了会堂？因为你是他的同党，你的分在来世与他同在。}\Person{尼苛德摩}说：\Quote{阿们，阿们。}同样，\Person{若瑟}也站出来对他们说：\Quote{你们为什么因为我求了耶稣的遗体而向我发怒？看，我把他放在我的新坟墓里，用洁净的细麻布包裹，又把一块石头滚到坟墓门口。你们对义人行了不善的事，因为你们没有悔改钉死他，还用长矛刺透了他。}犹太人抓住了\Person{若瑟}，下令将他拘禁到一周的第一天，并对他说：\Quote{要知道，时间不允许我们对你们做什么，因为安息日临近了；要知道，你不配得到埋葬，我们必将你的肉给空中的飞鸟。}\Person{若瑟}对他们说：\Quote{这是骄傲的\Person{哥肋雅}的话，他曾辱骂永生的天主和圣\Person{达味}。\BibleRef{撒上 17:44}因为天主曾藉先知说：复仇是我的，我必报应——上主说。如今那在肉体上未受割礼，心却受了割礼的人，取了水，在太阳面前洗手，说：\InnerQuote{我对这义人的血是无辜的；你们自己看着办吧。}你们却回答\Person{比拉多}说：\InnerQuote{他的血归在我们和我们子孙身上。}现在我害怕天主的义怒降临到你们和你们子孙身上，正如你们所说的。}犹太人听了这些话，心里愤恨，抓住了\Person{若瑟}，把他锁进一间没有窗户的房间里；门口有守卫看守，他们封住了关押\Person{若瑟}的门。\par

安息日，会堂的首领、司祭和肋未人下令，所有人要在一周的第一天到会堂集合。清早起来，会堂里的全体群众商议该用什么方法杀死他。当公议会开庭时，他们下令极其屈辱地把他带来。开了门，却没有找到他。全体民众都惊讶不已，惊慌失措，因为他们发现封条未破，钥匙也在\Person{盖法}手里。他们再也不敢向那些曾在\Person{比拉多}面前为\Person{耶稣}说话的人下手了。\par

% structure: p0058 source=h2 target=subsection reason=relative-heading-rank; base=h2

\subsection{第十三章}

当他们还在会堂里坐着，为\Person{若瑟}的事感到惊奇时，有几个守卫来了，就是犹太人曾向\Person{比拉多}请求来看守\Person{耶稣}坟墓的那几个，免得祂的门徒来偷走祂。他们向会堂的首领、司祭和肋未人报告了所发生的事：\Quote{如何有大地震；我们看见一位天使从天降下，把石头从坟墓门口滚开，坐在上面；他光芒如雪，又如闪电。我们非常害怕，像死人一样倒在地上；我们听见天使对守在坟墓旁的妇女说：\InnerQuote{不要害怕，我知道你们寻找被钉死的耶稣。祂不在这里：祂已复活，正如祂所说的。你们来看主躺过的地方；快去告诉祂的门徒，祂已从死者中复活，现在在\Place{加利利}。}\BibleRef{玛 28:5}}\par

犹太人说：\Quote{他对哪些妇女说了话？}守卫说：\Quote{我们不知道她们是谁。}犹太人说：\Quote{是什么时候的事？}守卫说：\Quote{半夜。}犹太人说：\Quote{你们为什么没有抓住她们？}守卫说：\Quote{我们因害怕像死人一样，不指望看见天光，又怎能抓住她们呢？}犹太人说：\Quote{上主永在，我们不信你们。}守卫对犹太人说：\Quote{你们在这人身上看见了这么大的奇迹，还是不信；又怎能相信我们？你们确实说了好话，指着永生的上主起誓，因为他的确活着。}守卫又说：\Quote{我们听说你们把求耶稣遗体的人关了起来，在门上加了封条；你们开了门，却没有找到他。那么，你们把你们看守的人交给我们，我们就把耶稣交给你们。}犹太人说：\Quote{\Person{若瑟}已经回他自己的城去了。}守卫对犹太人说：\Quote{\Person{耶稣}已经复活了，正如我们从天使那里听到的，现在在\Place{加利利}。}\par

犹太人听了这些话，非常害怕，说：\Quote{我们必须小心，免得这事传开，大家都归向\Person{耶稣}。}于是犹太人召集会议，付了一大笔钱，交给士兵，说：\Quote{你们要说，我们睡觉的时候，他的门徒夜间来把他偷走了；如果这事传到总督耳中，我们会说服他，让你们免于麻烦。}士兵拿了钱，就照所吩咐的说了。\par

% structure: p0062 source=h2 target=subsection reason=relative-heading-rank; base=h2

\subsection{第十四章}

\Person{斐尼阿}（司祭）、\Person{阿达}（教师）和\Person{哈盖}（肋未人）从\Place{加利利}下到\Place{耶路撒冷}，对会堂的首领、司祭和肋未人说：\Quote{我们看见\Person{耶稣}和祂的门徒坐在名叫\Place{玛米耳赫}的山上；祂对门徒说：\InnerQuote{你们往普天下去，向一切受造物宣讲福音；信而受洗的，必得救；不信的，必被定罪。信的人必有这些奇迹相随：因我的名驱逐魔鬼，说新语言，手拿毒蛇，甚至喝了致命的毒物也绝不受害，按手在病人身上，病人就会痊愈。}耶稣正对门徒说话的时候，我们看见祂被提升天。}\BibleRef{谷 16:15}\par

长老、司祭和肋未人说：\Quote{你们要光荣以色列的天主，并向祂承认，你们是否听见并看见了你们所告诉我们的事。}那报告的人说：\Quote{主永在，我们祖先亚巴郎、依撒格和雅各伯的天主，我们听见了这些事，并看见他被接升天。}长老、司祭和肋未人对他们说：\Quote{你们是来向我们报告这事，还是来向天主祈祷？}他们说：\Quote{是来向天主祈祷。}长老、司祭长和肋未人对他们说：\Quote{如果你们是来向天主祈祷，为什么在众人面前讲这些闲话？}司祭丕乃哈斯、教师阿达斯和肋未人哈盖对会堂长、司祭和肋未人说：\Quote{如果我们所说所见的罪大恶极，看，我们在你们面前，你们看怎样好，就怎样待我们吧。}他们便拿起法律书，叫他们指着它起誓，不再向任何人提及这些事。然后他们给他们吃喝，送他们出城，又给了他们钱和三个人作伴，送他们往加里肋亚去。这些人去了加里肋亚，司祭长、会堂长和长老都聚集在会堂里，锁上门，大哭一场说：\Quote{在以色列竟发生了这样的奇事？}亚纳斯和盖法说：\Quote{你们为什么如此激动？为什么哭泣？难道你们不知道，他的门徒给了守卫坟墓的人一笔金子，吩咐他们说，有天使下来，把石头从墓门口滚开了？}司祭和长老说：\Quote{就算他的门徒偷走了他的身体，但生命如何进入他的身体，他又如何在加里肋亚走动？}他们无法回答这些话，犹豫了很久说：\Quote{我们不可相信未受割礼的人。}\par

% structure: p0065 source=h2 target=subsection reason=relative-heading-rank; base=h2

\subsection{第十五章}

尼苛德摩站起来，站在公议会前说：\Quote{你们说得对，主的子民，你们并非不知道这些从加里肋亚下来的人，他们敬畏天主，是殷实之人，憎恨贪财，爱好和平；他们曾起誓说：\InnerQuote{我们看见耶稣在玛米耳克山上与门徒在一起，他教导了我们从他听到的话，并看见他被接升天。}没有人问他们他是以什么形式升天的。因为确实，正如圣经教训我们的，厄里亚也被接升天，厄里叟大声呼喊，厄里亚把羊皮抛在厄里叟身上，厄里叟把羊皮抛在约但河上，渡了过去，来到耶里哥。先知弟子们遇见他，说：\InnerQuote{厄里叟，你的师父厄里亚在哪里？}他说：\InnerQuote{他被接升天了。}他们对厄里叟说：\InnerQuote{莫非有灵抓住了他，把他扔在一座山上？让我们带着仆人去找他。}他们说服了厄里叟，他便同他们去了。他们找了三天，没有找到他，就知道他被接升天了。\BibleRef{列下 2:12} 现在请听我说，让我们打发人到以色列全境，看看基督是否也被灵抓住，扔在一座山上？}这个提议使众人喜悦。他们便打发人到以色列全境寻找耶稣，却没有找到他；但他们在阿黎玛特雅找到了若瑟，却没有人敢下手拿他。\par

他们向长老、司祭和肋未人报告说：\Quote{我们走遍了以色列全境，没有找到耶稣；但我们在阿黎玛特雅找到了若瑟。}他们听到若瑟的消息，十分高兴，就光荣以色列的天主。会堂长、司祭和肋未人商议如何与若瑟见面，便拿了一张纸，写给若瑟如下：—\par

\Quote{愿你平安！我们知道我们得罪了天主，也得罪了您；我们曾向以色列的天主祈祷，求您屈尊回到您的先辈和您的儿女那里，因为我们全都忧伤。因为我们开了门，没有找到您。我们知道我们曾谋害了您，但主保护了您，主亲自将我们对您的计谋吹散，可敬的若瑟父。}\par

他们从以色列全境选了七个人，都是若瑟的朋友，若瑟本人也认识他们；会堂长、司祭和肋未人对他们说：\Quote{你们注意：若他接到我们的信后读了，你们就知道他会同你们来见我们；但若他不读，就知道他对我们怀有恶意。你们向他请安后，就回到我们这里。}他们祝福了这些人，就打发他们走了。这些人来到若瑟那里，向他行礼，对他说：\Quote{愿你平安！}他说：\Quote{愿你们平安，也愿以色列全体人民平安！}他们便把信卷交给他。若瑟接过来，读了信，把它卷起来，祝福了天主说：\Quote{愿主天主受赞美，因祂拯救了以色列，使他们不流无辜的血；愿主受赞美，因祂派遣了祂的天使，用翅膀庇护了我。}他为他们摆设了筵席，他们便吃喝，并在那里过夜。清早他们起来祈祷。若瑟备好驴，同那些人出发，来到圣城耶路撒冷。全体百姓迎接若瑟，喊说：\Quote{愿你平安进入！}他对全体百姓说：\Quote{愿你们平安！}并亲吻了他们。百姓同若瑟祈祷，看见他都惊讶不已。尼苛德摩接他到自己家里，大摆筵席，请了亚纳斯、盖法、长老、司祭和肋未人来家里。他们同若瑟一起欢宴吃喝；唱完赞美诗后，各人回了自己的家。但若瑟留在尼苛德摩家里。\par

次日，即预备日，会堂长、司铎和肋未人清早来到\Person{尼苛德摩}的家；\Person{尼苛德摩}迎接他们，说：\Quote{愿你们平安！}他们对他说：\Quote{愿你们平安，也愿\Person{若瑟}平安，愿你的全家平安，也愿\Person{若瑟}的全家平安！}他便领他们进了自己的家。全体公议会都坐下了，\Person{若瑟}坐在\Person{亚纳斯}与\Person{盖法}之间；没有人敢对他说一句话。\Person{若瑟}说：\Quote{你们为什么召叫我来？}他们示意\Person{尼苛德摩}对\Person{若瑟}说话。\Person{尼苛德摩}开口对\Person{若瑟}说：\Quote{父亲，你知道那尊贵的经师、司铎和肋未人，要听你一句话。}\Person{若瑟}说：\Quote{请问吧。}\Person{亚纳斯}和\Person{盖法}拿起法律书，使\Person{若瑟}起誓说：\Quote{归光荣于以色列的天主，并向祂告明；因为\Person{阿干}被先知\Person{耶稣}起誓时，没有发虚誓，反而向他陈述了一切，没有向他隐藏一句话。你也照样不要向我们隐藏任何一句话。}\Person{若瑟}说：\Quote{我一句话也不向你们隐藏。}他们对他说：\Quote{我们非常忧伤，因为你求了\Person{耶稣}的身体，用洁净的殓布包裹，安放在坟墓里。因此，我们将你锁在一个没有窗户的房间，在门上加了锁和封条，又有守卫看守着你被锁之处。一周的第一天，我们开了门，却找不到你，极其忧伤；惊骇落在主的全体百姓身上，直到昨天。现在请告诉我们，你遭遇了什么事。}\par

\Person{若瑟}说：\Quote{在预备日，大约第十时辰，你们把我锁起来，我整个安息日都留在那里。午夜时分，我正站着祈祷，你们锁我的房间被四角吊起，我看见一道如闪电的光射入我的双眼。\BibleRef{宗 10:11}我害怕起来，跌倒在地。有人拉住我的手，把我从跌倒的地方扶起；水气从我的头浇到脚，一股香气扑鼻而来。他擦了我的脸，亲吻我，对我说：\InnerQuote{不要害怕，\Person{若瑟}；睁开你的眼睛，看看是谁在对你说话。}我抬头一看，看见了\Person{耶稣}。我战栗起来，以为是个幻影；我就背诵诫命，他也与我同诵。当然你们知道，幻影若遇见人又听见诫命，就会逃跑。我见他与我同诵，就对他说：\InnerQuote{辣彼\Person{厄里亚}。}他对我说：\InnerQuote{我不是\Person{厄里亚}。}我对他说：\InnerQuote{我主，你是谁？}他对我说：\InnerQuote{我是\Person{耶稣}，就是你向\Person{比拉多}求来身体的那位；你用洁净的殓布包裹我，把汗巾放在我脸上，将我安放在你的新坟墓里，又滚了一块大石头到墓门口。}我对那对我说话的说：\InnerQuote{请指给我看，我将你安放的地方。}他带我走，指给我看他被安放之处；殓布和脸上的汗巾都放在那里。我就知道那是\Person{耶稣}。他拉着我的手，虽然门是锁着的，却把我安置在我房子的中央，领我到我的床前，对我说：\InnerQuote{愿你平安！}他亲吻我，又对我说：\InnerQuote{四十天不要出门；因为，看，我往加里肋亚到我弟兄那里去。}}\par

% structure: p0072 source=h2 target=subsection reason=relative-heading-rank; base=h2

\subsection{第十六章}

会堂长、司铎和肋未人听了\Person{若瑟}这些话，如同死人，倒在地上，禁食到第九时辰。\Person{尼苛德摩}与\Person{若瑟}一同劝勉\Person{亚纳斯}、\Person{盖法}、司铎和肋未人说：\Quote{起来，站稳，尝点饼，坚固你们的灵魂，因为明天是主的安息日。}他们便起来，向天主祈祷，吃了喝了，各人回了自己的家。\par

安息日那天，我们经师、司铎和肋未人坐着彼此询问说：\Quote{这临到我们的愤怒是什么？因为我們知道他的父母。}\Person{肋未}，一位经师，说：\Quote{我知道他的父母敬畏天主，不远离祈祷，一年三次献什一之物。\Person{耶稣}诞生时，他的父母带他来到这地方，向天主献了祭献和全燔祭。伟大的经师\Person{西默盎}把他抱在怀里，说：\BibleQuote{主啊！现在祢按照祢的话，放祢的仆人平安去吧；因为我的眼睛已看见祢的救恩，就是祢在万民面前所预备的：为启示外邦人的光，祢的百姓以色列的荣耀。}\Person{西默盎}祝福了他们，对他的母亲\Person{玛利亚}说：\InnerQuote{关于这孩子，我给你们报告好消息。}\Person{玛利亚}说：\InnerQuote{好，我主。}\Person{西默盎}对她说：\InnerQuote{好；}\BibleQuote{看，这孩子成了许多人在以色列中跌倒和兴起的原因，并成为反对的标记；至于你，一把剑将刺透你自己的灵魂，为叫许多人的心思意念得以显露。}} \BibleRef{路 2:25}\par

他们对经师\Person{肋未}说：\Quote{你怎么知道这些事？}\Person{肋未}对他们说：\Quote{难道你们不知道我从他学了法律？}公议会问他：\Quote{我们愿见你的父亲。}他们便派人去叫他的父亲。他们问了他；他对他们说：\Quote{你们为什么不信我的儿子？那位有福且正义的\Person{西默盎}亲自教了他法律。}公议会问辣彼\Person{肋未}：\Quote{你所说的话是真的吗？}他答说：\Quote{是真的。}会堂长、司铎和肋未人彼此说：\Quote{来，我们派人往加里肋亚，到那三个曾来告诉我们关于他的教导和被接升天的人那里，让他们告诉我们他们如何看见他被接升天。}这话使众人喜悦。他们派走了那三个已经同他们往加里肋亚去的人；对他们说：\Quote{对辣彼\Person{阿达斯}、辣彼\Person{丕乃哈斯}和辣彼\Person{哈盖}说：愿你们平安，也愿一切与你们同在的人平安！在公议会中发生了重大查问，我们被派到你们这里，召你们到圣地耶路撒冷。}\par

那些人便动身往\Place{加里肋亚}去，找到他们坐着研读法律，就平安地向他们请安。在\Place{加里肋亚}的那些人对来到他们这里的人说：\Quote{愿平安归于全以色列！}他们回答说：\Quote{平安归于你们！}他们又问他们说：\Quote{你们为什么来？}被派去的人说：\Quote{公议会召你们到圣城\Place{耶路撒冷}。}那些人听见自己被公议会寻找，就向天主祈祷，然后同那些人一起斜倚用餐，吃了喝了，就起身平安地往\Place{耶路撒冷}去了。\par

次日，公议会坐在会堂中，问他们说：\Quote{你们果真看见\Person{耶稣}坐在\Place{玛米耳山}上教导他的十一位门徒，又看见他被接升天吗？}那些人回答他们说：\Quote{我们怎样看见他被接升天，也就怎样说了。}\par

\Person{亚纳斯}说：\Quote{把他们彼此分开，让我们看看他们的陈述是否一致。}他们就把他们彼此分开。首先召来\Person{阿达斯}，对他说：\Quote{你怎样看见\Person{耶稣}被接升天？}\Person{阿达斯}说：\Quote{当他还在\Place{玛米耳山}上教导他的门徒时，我们看见一朵云彩遮盖了他和他的门徒。云彩把他接到天上去了，他的门徒都脸伏于地。}他们又召来司祭\Person{非乃哈斯}，也问他说：\Quote{你怎样看见\Person{耶稣}被接升天？}他照样说了。他们又问\Person{哈盖}，他也照样说了。公议会说：\Quote{\BibleQuote{凭两个或三个人的口，一切陈述都可成立。}}\BibleRef{申 17:6} 教师\Person{布特姆}说：\Quote{\BibleQuote{哈诺客与天主同行，就不见了，因为天主将他取去。}}教师\Person{雅依洛}说：\Quote{\BibleQuote{我们听说了圣梅瑟的死，却没有看见；因为上主的法律上记载：梅瑟因上主的话而死，直到今日没有人知道他的坟墓。}\BibleRef{申 34:5}} 辣彼\Person{肋未}说：\Quote{为什么辣彼\Person{西默盎}看见\Person{耶稣}时说：\BibleQuote{看，这孩子将是以色列中许多人跌倒和兴起的原因，并成为反对的记号？}\BibleRef{路 2:34}} 辣彼\Person{依撒格}说：\Quote{\BibleQuote{看，我派遣我的使者在你面前，他要在你前面行，在路上保护你，使你进入我所预备的地方，因为我的名在他身上。}}\par

那时\Person{亚纳斯}和\Person{盖法}说：\Quote{你们说得对，梅瑟法律上所记载的，没有人看见哈诺客的死，也没有人说起梅瑟的死；但\Person{耶稣}在\Person{比拉多}面前受审，我们看见他挨打，脸上被吐唾沫，兵士给他戴上茨冠，他被鞭打，受了\Person{比拉多}的判决，在\Place{髑髅地}被钉十字架，同他一起还有两个强盗；他们给他喝醋和苦胆调和的酒，兵士\Person{隆季诺}用长枪刺透了他的肋旁；我们可敬的父亲\Person{若瑟}求得了他的遗体，并且如他所说，他已经复活了；正如那三位教师所说，我们看见他被接升天；辣彼\Person{肋未}也证明了辣彼\Person{西默盎}所说的话，他说：\BibleQuote{看，这孩子将是以色列中许多人跌倒和兴起的原因，并成为反对的记号。}\BibleRef{路 2:34}} 所有的教师对上主的全体子民说：\Quote{这事若是出于上主，在你们眼中看为奇妙，你们就该知道，雅各伯家啊，法律上记载：\BibleQuote{凡挂在树上的，都是可咒骂的。} 另一段圣经教训说：\BibleQuote{那些没有创造天地的神祇，必从地上、从天下被除灭。}\BibleRef{耶 10:11}} 司祭和肋未人彼此说：\Quote{如果他的纪念直到称为禧年的那一年，你们要知道，它将永远存留，并且他为自己兴起了一个新子民。}于是会堂的领袖、司祭和肋未人向全以色列宣布说：\Quote{凡敬拜人手所造之物的，是可咒骂的；凡敬拜受造物超过造物主的，也是可咒骂的。}全体人民说：\Quote{阿们，阿们。}\par

全体人民赞美上主说：\BibleQuote{上主是应当称颂的，因为他照着他所说的一切，赐给了他的子民以色列安息；他的一切善言，没有一句落空，就是他藉他的仆人梅瑟所说的话。愿上主我们的天主与我们同在，如同与我们的祖先同在一样；愿他不毁灭我们。求他不毁灭我们，使我们专心归向他，遵行他的一切道路，遵守他吩咐我们祖先的诫命和典章。}\BibleRef{列上 8:56} \BibleQuote{那日，上主必为普世的君王；只有一位上主，他的名是唯一的。}\BibleRef{匝 14:9} \BibleQuote{上主是我们的君王，他必拯救我们。}\BibleRef{依 33:22} \BibleQuote{上主啊，没有谁能与你相比。上主啊，你是伟大的，你的名也是伟大的。求你用你的大能医治我们，我们就痊愈；拯救我们，我们就得救。}\BibleRef{耶 17:14} \BibleQuote{因为我们是你的产业和基业。上主为自己的大名，必不丢弃他的子民；因为上主已开始使我们成为他的子民。}\BibleRef{撒上 12:22}\par

众人唱完赞美诗，就各自回家去了，光荣天主，因为光荣属于他，直到永远。阿们。\par

\SourceRecord{Translated by Alexander Walker.；Ante-Nicene Fathers；Vol. 8.；Edited by Alexander Roberts, James Donaldson, and A. Cleveland Coxe.；Buffalo, NY: Christian Literature Publishing Co.,；1886.；<http://www.newadvent.org/fathers/08071a.htm>.}

% structure: p0001 source=h1 target=section reason=page-title

\section{尼苛德摩福音（第一部分，第二种希腊文本）}

叙述关于我们的主耶稣基督的苦难和祂的圣复活。\par

由一个名叫厄奈阿斯的犹太人撰写，并由一位罗马地方长官尼苛德摩从希伯来文翻译成罗马语。\par

在希伯来王国解体后，四百年流逝，希伯来人最终也置于罗马帝国统治之下，罗马王为他们立了一位王；当提庇留凯撒执掌罗马权杖时，在他统治的第十八年，他立了黑落德——即从前在白冷屠杀婴儿的那位黑落德之子——为犹太王，并任命比拉多为耶路撒冷总督。当亚纳斯和盖法担任耶路撒冷的大司祭时，罗马地方长官尼苛德摩召来一个名叫厄奈阿斯的犹太人，请他撰写一份关于亚纳斯和盖法时代在耶路撒冷所行关于基督之事的记录。这犹太人便照做了，并将记录交给尼苛德摩；尼苛德摩又将其从希伯来文翻译成罗马语。记录如下：\par

% structure: p0005 source=h2 target=subsection reason=relative-heading-rank; base=h2

\subsection{第一章}

我们的主耶稣基督在犹太行了许多伟大而奇妙的奇迹，并因此被希伯来人仇恨。当比拉多做耶路撒冷总督，亚纳斯和盖法做大司祭时，有些犹太人来到司祭长面前，有犹达斯、肋未、纳斐塔里、亚历山大、西罗及其他许多人，他们指责基督。这些司祭长便派他们去对比拉多也讲这些事。他们就去对他说：在这城里行走着一个人，他的父亲叫若瑟，母亲叫玛利亚；他自称是君王和天主子；他身为犹太人，却推翻圣经，废除安息日。比拉多于是询问，想从他们口中得知他是如何废除安息日的。他们回答说：他在安息日治愈病人。比拉多说：若他使病人痊愈，他并没有作恶。他们对他说：若他妥善地治病，恶尚小；但他却用魔术行这些事，并依靠魔鬼。比拉多说：治愈病人不是魔鬼的作为，而是来自天主的恩宠。\par

希伯来人说：我们恳请阁下召见他，以便您能确切查问我们所说的话。于是比拉多脱下外袍，交给一个差役，说：去，把这件外袍拿给耶稣看，并对他说，比拉多总督召你到他面前来。差役便去了，找到耶稣，召叫他，同时将比拉多的外袍铺在地上，并请耶稣踏在上面。希伯来人见此情景，极为愤怒，来到比拉多面前，埋怨他竟赐予耶稣如此大的尊荣。\par

比拉多询问被派去的差役为何如此行，差役回答说：当您派我到犹太人亚历山大那里时，我遇见耶稣正骑着一头驴进入城门。我看见希伯来人将衣服铺在路上，驴就从衣服上踏过；另有人砍下树枝，前去迎接他，并高喊：贺三纳于高天！因此，我也必须这样做。\par

犹太人听了这些话，对他说：你身为罗马人，如何知道希伯来人所说的话？差役回答说：我问了一个希伯来人，是他告诉我的。比拉多说：贺三纳是什么意思？犹太人说：主啊，求祢拯救。比拉多回答说：既然你们承认你们的子女这样说了，如今你们又怎会指控耶稣，说出你们所说的那些话呢？犹太人沉默了，无话可答。\par

当耶稣来到比拉多面前时，比拉多的士兵朝拜了他。还有其他人手持旗帜站在比拉多面前。当耶稣进来时，旗帜也弯下腰，朝拜了他。比拉多对所见之事感到惊讶，犹太人对他说：我主，并非旗帜朝拜耶稣，而是举旗的士兵随意所致。\par

比拉多对会堂长说：挑选十二个强壮的人，把旗帜交给他们，让他们牢牢握住。这事发生后，比拉多命令差役将耶稣带到外面，再带进来。当耶稣进来时，旗帜再次弯下腰，朝拜了他。比拉多因此大为惊异。但犹太人说：他是个术士，藉此行这些事。\par

% structure: p0012 source=h2 target=subsection reason=relative-heading-rank; base=h2

\subsection{第二章}

比拉多对耶稣说：祢听见这些人作证反对祢吗？祢为什么不回答？\BibleRef{玛 27:13}耶稣回答说：每个人都有能力说好或说坏，随他所愿；因此这些人也有能力，说他们愿意说的。\BibleRef{若 19:11}\par

犹太人对他说：关于祢我们该说什么？首先，祢是从罪而生的；其次，因为祢，当祢出生时，婴儿被杀害；第三，祢的父亲和母亲逃往埃及，因为他们对人民没有信心。\par

对此，当时在场的敬畏天主的犹太人回答说：我们说他的出生并非来自罪恶；因为我们知道若瑟是按照订婚习俗接受了他的母亲玛利亚。比拉多说：因此你们说他的出生来自罪恶，是在说谎。他们又对比拉多说：全体民众都证明他是个术士。敬畏天主的犹太人回答说：我们也在他母亲的订婚现场，我们是犹太人，了解他每日的生活；但他是否是术士，我们不知道。这样说的犹太人有：拉匝禄、阿斯达里、安东尼、雅各伯、匝辣、撒慕尔、依撒格、斐尼、克里斯普、达格里普、阿默色、犹达斯。\par

于是比拉多对他们说：凭凯撒的性命，我希望你们起誓，这个人的出生是否无罪。他们回答说：我们的法律规定我们绝不可起誓，因为起誓是大罪。然而，凭凯撒的性命我们起誓，他的出生是无罪的；如果我们说谎，请下令将我们全部斩首。他们这样说了之后，提出指控的犹太人回答比拉多说：难道你相信这十二个犹太人，胜过相信全体民众和我们吗？我们确切知道他是术士、亵渎者，且自称为天主子。\par

于是彼拉多命令所有人退出总督府，只留下那十二人。此事既成，彼拉多私下对他们说：\Quote{关于这人，看来犹太人出于嫉妒和疯狂想杀他：因为他们控告他做了一件事——他废止安息日；但他那时却做了善工，因为他治愈病人。为此，这人不应被判死刑。}那十二人也对他说：\Quote{确实，我主，正是如此。}\par

% structure: p0018 source=h2 target=subsection reason=relative-heading-rank; base=h2

\subsection{第三章}

因此彼拉多愤慨且恼怒地走到外面，对亚纳斯和盖法以及带来耶稣的群众说：\Quote{我以太阳为证，我在此人身上找不到任何罪过。}群众回答说：\Quote{如果他不是巫士、术士和亵渎者，我们就不会把他带到阁下面前。}彼拉多说：\Quote{你们自己审问他；既然你们有法律，就按你们的法律办。}犹太人说：\Quote{我们的法律不准处死人。} \BibleRef{若 19:6}彼拉多说：\Quote{如果你们不愿意处死他，那我更不愿意！}\par

然后彼拉多回到府内，对耶稣说：\Quote{告诉我，你是犹太人的王吗？}耶稣回答说：\Quote{你自己说这话，还是别的犹太人告诉你的，好让你来问我？}彼拉多说：\Quote{你以为我是希伯来人吗？我不是希伯来人。你的百姓和司祭长把你交在我手里；告诉我你是不是犹太人的王？}耶稣回答说：\Quote{我的国不属于这世界；如果我的国属于这世界，我的士兵就不会对我被捕无动于衷：因此我的国不在这世界上。}彼拉多说：\Quote{既然如此，你是王吗？}耶稣说：\Quote{你说的是：我为此而生，要为真理作证；凡属于真理的人，就听我的话，并实行它。}彼拉多说：\Quote{什么是真理？} \BibleRef{若 18:33}耶稣回答说：\Quote{真理来自天上。}彼拉多说：\Quote{那么，在地上就没有真理吗？}基督说：\Quote{我就是真理；那有地上权力的人如何审判真理呢！}\par

% structure: p0021 source=h2 target=subsection reason=relative-heading-rank; base=h2

\subsection{第四章}

因此彼拉多离开基督，走到外面，对犹太人说：\Quote{我在这人身上找不到任何罪过。}犹太人回答说：\Quote{让我们告诉阁下他说了什么。他说：\InnerQuote{我能拆毁天主的圣殿，并在三天内重建它。}}彼拉多说：\Quote{他说要拆毁哪座圣殿？}希伯来人说：\Quote{所罗门的圣殿，所罗门用了四十六年建成的。} \BibleRef{若 2:20}\par

彼拉多私下对司祭长、经师和法利塞人说：\Quote{我恳求你们，不要对这人行恶；因为你们若对他行恶，你们就不义：这样一个人不应死，他给许多人行了极大的善事。}他们对彼拉多说：\Quote{我主，若是侮辱凯撒的人该死，那么这侮辱天主的\Emph{人}岂不更该死！}\par

于是彼拉多遣散他们，他们都出去了。随后他对耶稣说：\Quote{你希望我怎样处置你？}耶稣对彼拉多说：\Quote{照所定的办吧。}彼拉多说：\Quote{如何定的？}耶稣回答说：\Quote{梅瑟和先知们为我被钉十字架并复活的事写过。}希伯来人听了\Emph{这话}，对彼拉多说：\Quote{你为什么还要听他更大对天主的侮辱？}彼拉多说：\Quote{这些话不是对天主的侮辱，因为它们写在先知书上。}希伯来人说：\Quote{我们的圣经说：若有人得罪人，即侮辱人，他应受四十杖；但若有人侮辱天主，则应被石头砸死。}\par

那时，彼拉多的妻子普罗克拉派一个使者到他那里；信上说：\Quote{你要当心，不要同意让这义人耶稣受害；因为今夜我因他而作了可怕的梦。} \BibleRef{玛 27:19}彼拉多对希伯来人说：\Quote{如果你们认为耶稣所说的话是对天主的侮辱，就按你们的法律自己审问他吧。} \BibleRef{若 18:31}犹太人对彼拉多说：\Quote{我们愿你把他钉十字架。}彼拉多说：\Quote{这不好。}\par

彼拉多转身看见许多人哭泣，就说：\Quote{在我看来，不是所有人都希望这人死。}司祭和经师说：\Quote{我们为此带来了所有人，好让你完全确信所有人都希望他死。}彼拉多说：\Quote{他做了什么恶事？}希伯来人说：\Quote{他说他是王，是天主子。}\par

% structure: p0027 source=h2 target=subsection reason=relative-heading-rank; base=h2

\subsection{第五章}

于是有一个敬畏天主的犹太人，名叫尼苛德摩，站出来对彼拉多说：\Quote{我恳请阁下允许我说几句话。}彼拉多说：\Quote{请说吧。}尼苛德摩说：\Quote{我曾在会堂中对司祭、肋未人、经师和百姓说：你们要怎样反对这人？这人行了许多神迹，是前人未曾行过、也未曾行过的。所以，放了他吧；如果他做的事是从天主来的，必要存立；若是从人来的，必要消灭。} \BibleRef{宗 5:38} \Quote{正如天主派遣梅瑟到埃及时所发生的一样：埃及王法郎叫他行神迹，他就行了。那时法郎有两个术士，雅乃和杨布雷；他们也用巫术行了神迹，却不像梅瑟所行的。} \BibleRef{出 7:10} \Quote{埃及人把这些术士当作神；但因他们不是从天主来的，他们所做的就被消灭了。这耶稣曾使拉匝禄复活，他还活着。为此我恳求你，我主，绝不可让这人被杀。}\par

希伯来人对尼苛德摩发怒，说：\Quote{愿你接受耶稣的真理，并与他同分。}尼苛德摩说：\Quote{阿们，阿们；愿照你们所说的成就于我。}\par

% structure: p0030 source=h2 target=subsection reason=relative-heading-rank; base=h2

\subsection{第六章}

尼苛德摩说完后，另一个希伯来人站起来对彼拉多说：\Quote{我求你，我主彼拉多，也请听我说。}彼拉多回答说：\Quote{说你想说的。}那个希伯来人说：\Quote{我患病卧床三十八年；他看见我就动了怜悯之心，对我说：\InnerQuote{起来，拿起你的床，回家去吧。}他说这话时，我就起来行走了。}希伯来人说：\Quote{问他这事发生在哪一天。}他说：\Quote{在安息日。} \BibleRef{若 5:5}犹太人说：\Quote{因此我们确实说，他不守安息日。}\par

另有一人立于人群中，说：\Quote{我生来就瞎；耶稣沿路行时，我向他呼喊说：\InnerQuote{主，达味之子，可怜我吧！}他取了泥，抹在我的眼上，我立刻就看见了。}\BibleRef{若 9:6}又有一人说：\Quote{我生来驼背；我看见他，就呼喊：\InnerQuote{主，可怜我吧！}他抓住我的手，我立刻就被扶直了。}\BibleRef{宗 3:7}又有一人说：\Quote{我曾是癞病人，他只凭一句话就治愈了我。}\BibleRef{路 17:11}\par

% structure: p0033 source=h2 target=subsection reason=relative-heading-rank; base=h2

\subsection{第七章}

有一个名叫\Person{韦罗尼加}的女人也在那里，她说：\Quote{我患血漏十二年，我只摸了摸他衣边的穗子，立刻就痊愈了。}\BibleRef{玛 9:20}犹太人说：\Quote{我们的法律不接受女人的证言。}\par

% structure: p0035 source=h2 target=subsection reason=relative-heading-rank; base=h2

\subsection{第八章}

另有人喊道：\Quote{这人是先知，连魔鬼都怕他。}\Person{比拉多}说：\Quote{那么，魔鬼为何从未这样怕过你们的祖先呢？}他们说：\Quote{我们不知道。}又有人说：\Quote{\Person{拉匝禄}在坟墓里已有四天，他只凭一句话就使他复活了。}\BibleRef{若 11:43}\Person{比拉多}听了\Person{拉匝禄}复活的事，就害怕，对百姓说：\Quote{你们为什么想要流这义人的血呢？}\par

% structure: p0037 source=h2 target=subsection reason=relative-heading-rank; base=h2

\subsection{第九章}

于是他召来\Person{尼苛德摩}和十二个敬畏天主的犹太人，对他们说：\Quote{你们说我该怎么做？因为百姓在骚动。}他们说：\Quote{我们不知道；你愿意怎样就怎样；但百姓所做的，是为了杀害他而行不义。}\Person{比拉多}又出去对百姓说：\Quote{你们知道，在无酵节期间，我习惯为你们释放一个在押的囚犯。我现在狱中有一个罪犯，是强盗，名叫\Person{巴拉巴}。我也有耶稣，他从未作过任何恶事。那么，你们愿意我把哪一个释放给你们？}百姓回答：\Quote{把\Person{巴拉巴}释放给我们。}\Person{比拉多}说：\Quote{那么，我拿耶稣怎么办？}他们说：\Quote{把他钉在十字架上。}又有人喊道：\Quote{你若释放耶稣，你就不是\Person{凯撒}的朋友，\BibleRef{若 19:12}因为他自称是天主子，是君王。你若释放他，他就成了君王，要夺去\Person{凯撒}的王国。}\par

\Person{比拉多}于是发怒，说：\Quote{你们的民族总是魔鬼一般，不信不忠，你们始终与你们的恩主为敌。}希伯来人说：\Quote{谁是我们的恩主？}\Person{比拉多}说：\Quote{天主，他从\Person{法老}手中解救你们，带领你们过红海如履旱地，赐你们鹌鹑为食，从干石中赐水喝，又赐给你们法律，你们却背弃天主，违背法律；若不是\Person{梅瑟}站立并恳求天主，你们早已死于惨痛的死亡。这一切你们都忘记了。同样，如今你们竟说我不爱\Person{凯撒}，反而恨他，想要阴谋反对他的王国。}\par

说完这些，\Person{比拉多}怒气冲冲地从宝座上站起来，想要离开他们。于是犹太人大喊说：\Quote{我们愿\Person{凯撒}做我们的王，不要耶稣，因为耶稣曾从贤士那里接受礼物。}连\Person{黑落德}也听说——将要有一位君王——想要杀死他，为此他派人杀尽了\Place{白冷}所有的婴儿。正因如此，他的父亲\Person{若瑟}和母亲因怕\Person{黑落德}，逃往\Place{埃及}。\par

\Person{比拉多}听了这话，就叫众人安静，说：\Quote{那么，这个耶稣就是\Person{黑落德}那时想要杀害的？}他们对他说：\Quote{是的。}\Person{比拉多}查明耶稣属\Person{黑落德}管辖，因他是犹太人血统，就把他送到\Person{黑落德}那里。\Person{黑落德}看见耶稣，非常欢喜，因为他早就想要见他，听说了他所行的奇迹。于是他给他穿上白袍，然后开始问他。但耶稣没有回答他。\Person{黑落德}想要看耶稣行什么奇迹，却没有看到，又因耶稣没有回答他一句话，就把他送回\Person{比拉多}那里。\Person{比拉多}看见这事，就命令官兵拿水来。他用水洗手，对百姓说：\Quote{对这义人的血，我是无罪的。你们要留意，他是不义地被处死，因为我和\Person{黑落德}都没有查出他有什么罪；正因如此，他才把他送回我这里。}犹太人说：\BibleQuote{他的血归在我们和我们子孙的身上！}\BibleRef{玛 27:25}\par

于是\Person{比拉多}坐在审判席上宣判。他下令，让耶稣来到他面前。他们拿了一个荆棘冠冕，戴在他头上，把一根芦苇放在他右手里。然后他宣判，对他说：\Quote{你的百姓指控你，证明你想要做王。因此我判决：按王国的法律，先用杖打你四十下，然后戏弄你，最后把你钉在十字架上。}\par

% structure: p0043 source=h2 target=subsection reason=relative-heading-rank; base=h2

\subsection{第十章}

\Person{比拉多}这样宣判后，犹太人开始击打耶稣，有的用棍棒，有的用手，有的用脚；还有人向他脸上吐唾沫。他们立刻准备好十字架，交给他，急忙上路。他这样走着，背着十字架，来到\Place{耶路撒冷}城门口。但因他受了许多击打，加上十字架的重压，无法行走，犹太人急于尽快钉死他，就从他身上取下十字架，交给一个迎面而来的人，名叫\Person{西满}，他有两个儿子，\Person{亚历山大}和\Person{鲁富}。他是\Place{基勒乃}城的人。\BibleRef{谷 15:21}他们把十字架交给他，并非因为可怜耶稣，想减轻他的负担，而是因为他们急于（如前所述）尽快处死他。\par

因此，祂的门徒中，\Person{若望}跟随着祂。然后他跑来见天主之母，对她说：\Quote{你去了哪里，怎么没来看发生了什么事？}她回答说：\Quote{发生了什么事？}\Person{若望}说：\Quote{你要知道，犹太人已捉住我的师傅，要带去钉死祂。}祂的母亲听见这话，就大声呼喊说：\Quote{我的儿子，我的儿子，你究竟做了什么恶事，他们竟要带你去处死？}她好像瞎了眼似的站起来，沿着路哭着走。妇女们跟随着她——\Person{玛尔大}、\Person{玛利亚玛达肋纳}、\Person{撒罗默}，以及其他贞女。\Person{若望}也与她在一起。当他们来到人群聚集的地方时，天主之母对\Person{若望}说：\Quote{我的儿子在哪里？}\Person{若望}说：\Quote{你看见那戴着荆棘冠冕、双手被绑的祂吗？}天主之母听见这话并看见祂，就晕倒，仰面跌倒在地，躺了很长时间。那些跟随她的妇女都围着她哭泣。她一苏醒过来，就大声呼喊说：\Quote{我的主，我的儿子，你俊美的容貌哪里去了？我怎能忍受看见你受这样的苦？}说着，她用指甲抓脸，捶胸。她说：\Quote{你在犹太所做的善事到哪里去了？你对犹太人做了什么恶事？}犹太人看见她这样哀哭，就来把她从路上赶走；但她不肯逃走，仍留在那里，说：\Quote{你们这些无法无天的犹太人，先杀死我吧。}\par

然后他们安全地到达了一个叫\Place{髑髅}的地方，是用石头铺成的；犹太人就在那里竖起了十字架。然后他们剥了耶稣的衣服，士兵们拿了祂的衣裳，彼此分着；他们给祂穿上一件破旧的猩红袍子，把祂举起来，在第六时辰把祂挂在十字架上。之后，他们又带来了两个强盗，一个在祂右边，一个在祂左边。\par

那时，天主之母站着观看，大声呼喊说：\Quote{我的儿子！我的儿子！}耶稣转过身来，看见她旁边站着\Person{若望}，并与其他妇女一同哭泣，就说：\Quote{看，你的儿子！}然后又对\Person{若望}说：\Quote{看，你的母亲！}\BibleRef{若 19:26}她大哭说：\Quote{为此我哭泣，我的儿子，因为你无辜受苦，因为无法无天的犹太人把你交到惨死。没有你，我的儿子，我将如何？没有你，我怎么活？我要过怎样的生活？你的门徒们在哪里？他们曾夸口要与你同死。那些被你治愈的人在哪里？怎么没有人来帮助你？}她又望着十字架说：\Quote{十字架，请你弯下，好让我拥抱亲吻我的儿子，就是我用这些乳房奇妙地哺乳过的，因为我未曾认识男人。十字架，请你弯下，我想拥抱我的儿子。十字架，请你弯下，好让我像一个母亲向我的儿子告别。}犹太人听见这些话，就上前来，把她和妇女们及\Person{若望}都赶到了远处。\par

那时耶稣大声呼喊说：\Quote{父啊，不要将这罪归于他们，因为他们不知道自己在做什么。}然后祂说：\Quote{我渴。}立刻有一个士兵跑去，拿了一块海绵，蘸满了苦胆和醋混成的液体，放在一根芦苇上，递给耶稣喝。祂尝了后，不肯喝。犹太人在一旁站着观看，嘲笑祂说：\Quote{你果真说你是天主子，就从十字架上下来，立刻，让我们相信你。}另一些人嘲弄说：\Quote{他救了别人，他治好了别人，他医好了病人、瘫痪者、癞病人、附魔者、瞎子、跛子、死人；自己却不能医治自己。}\BibleRef{玛 27:40}\par

同样，被钉在祂左边的强盗也对祂说：\Quote{如果你是天主子，就下来拯救你自己和我们。}他的名字叫\OriginalTerm{基斯塔}{Gistas}。被钉在右边的，名叫\OriginalTerm{狄斯玛}{Dysmas}，责备那强盗说：\Quote{你这可怜又可悲的人，难道你不敬畏天主吗？我们受的惩罚是我们应得的；但这个人没有做过任何恶事。}然后转向耶稣，对祂说：\Quote{主啊，当祢为王时，不要忘记我。}祂对他说：\Quote{今天，我实话告诉你，你将要与我一同在乐园里。}\par

% structure: p0050 source=h2 target=subsection reason=relative-heading-rank; base=h2

\subsection{第十一章}

那时耶稣大声呼喊说：\Quote{父啊，我将我的灵魂交在祢手中}，就断了气。\BibleRef{路 23:46}立刻可见磐石崩裂，因为全地发生地震；由于地震又强又大，磐石也都崩裂了。死者的坟墓也开了，圣所的帐幔裂为两半，从第六时辰到第九时辰遍地黑暗。犹太人看见这一切发生，就害怕，说：\Quote{这确实是个义人。}站在那里的百夫长\Person{隆基诺}说：\Quote{这真是天主子。}其他人来看见祂，因恐惧而捶胸，然后又回去了。\BibleRef{路 23:44}\par

百夫长看见了这一切如此大的奇迹，就去禀报给\Person{比拉多}。\Person{比拉多}听了，诧异又惊奇，因恐惧和悲伤，那日不吃不喝。他派人通知，黑暗过后，全体\OriginalTerm{公议会}{Sanhedrin}都来见他；他对民众说：\Quote{你们知道太阳如何变暗了；你们知道帐幔如何裂开了。我确实做得好，根本不愿处死这义人。}但恶徒们对比拉多说：\Quote{这黑暗是日食，如同其他时候发生过的一样。}然后他们对他说：\Quote{明天我们过无酵节；恳求你，既然被钉的人还在呼吸，打断他们的腿，把他们放下来。}\Person{比拉多}说：\Quote{就这样办。}于是他派兵去，发现两个强盗还在呼吸，就打断了他们的腿；但发现耶稣已死，他们丝毫未碰祂，只有一个士兵用枪刺穿了祂的右肋，立刻流出血和水来。\BibleRef{若 19:31}\par

\Emph{预备日}将近黄昏，若瑟——一个出身高贵、富有、敬畏天主的犹太人——找到尼苛德摩，后者先前所说的话已表明其心迹，便对他说：\Quote{我知道你爱耶稣在世时，乐意听祂的话，我也看见你为祂的缘故与犹太人争辩。因此，若你认为合宜，我们可到彼拉多那里，为埋葬而求耶稣的身体，因为祂未得安葬是一项大罪。}尼苛德摩说：\Quote{我怕彼拉多会发怒，灾祸临到我身上。但若你独自去求那死者并领祂，我便与你同去，帮你办理安葬所需的一切。}尼苛德摩这样说了，若瑟便举目向天，祈求不要使他的请求落空；于是他往彼拉多那里，向他请安后便坐下。然后他对彼拉多说：\Quote{我请求我主，不要向我发怒，倘若我所求的有悖于阁下心意。}彼拉多说：\Quote{你所求的是什么？}若瑟说：\Quote{耶稣——这善人，因犹太人的仇恨而被拿去钉死——我求你将他赐给我安葬。}彼拉多说：\Quote{这是怎么回事？我们竟要将那因本族送来巫术证据、且有篡夺凯撒王国嫌疑、因而被我们处死之人的尸体再予以尊荣？}若瑟流泪痛哭，俯伏在彼拉多脚前说：\Quote{我主，不要让仇恨落在死者身上；因为人所作的一切恶，应在死亡中与他同归于尽。我知道阁下曾何等迫切地不愿耶稣被钉死，也知你为祂的缘故向犹太人说了多少话——时而恳求，时而愤怒——最后你洗手表明自己决不与那些要处死祂的人同流合污。为这一切\Emph{理由}，我恳求你不要拒绝我的请求。}彼拉多见若瑟如此俯伏哀求哭泣，便扶起他说：\Quote{去吧，我将这死者赐给你；领去，随意处理。}\par

于是若瑟感谢了彼拉多，亲吻他的手和衣裳，便出去了。他心中固然为如愿以偿而喜乐，眼中却带着泪水。这样，他虽忧伤，却也欢喜。于是他往尼苛德摩那里去，将所发生的一切告诉他。然后他们买了一百斤的没药和沉香，并一座新坟墓\BibleRef{玛 27:60}，同天主之母、玛利亚玛达肋纳、撒罗默，以及若望和其余妇女，按惯例用白殓布处理了尸体，放在坟墓中。\BibleRef{若 19:38}\par

天主之母哭泣说：\Quote{我怎能不为你哀恸，我儿？我怎能不用指甲抓破我的脸？我儿，这就是西默盎长老在我抱你——四十天的婴孩——进圣殿时向我预言的。这剑如今刺透我的灵魂。\BibleRef{路 2:35}谁能止住我的眼泪，我最甘饴的儿？除你以外，谁也不能；除非你照你所说的，在三天后复活。}\par

玛利亚玛达肋纳哭泣说：\Quote{万民、万族、万舌啊，请听，要知道无法无天的犹太人将那位曾向他们行过无数善事者置于何等死地。请听，且要惊骇。谁能使这些事传遍天下？我要独自往罗马去见凯撒，向他指明彼拉多因顺从无法无天的犹太人而行了何等恶事。}同样，若瑟也哀悼说：\Quote{哀哉！最甘饴的耶稣，人中至善者——若可称你为人——你行了无人能行的奇事。我如何为你裹尸？如何将你安葬？如今那些你用几个饼养活的人应在场；这样，我便不至于显得未尽本分。}\par

然后若瑟与尼苛德摩回家去了；同样，天主之母与妇女们也回去了，若望也与他们同在。\par

% structure: p0058 source=h2 target=subsection reason=relative-heading-rank; base=h2

\subsection{第十二章}

犹太人得知若瑟和尼苛德摩所行的事，便对二人极其恼怒。司祭长亚纳斯和盖法派人去叫若瑟来，说：\Quote{你为什么为耶稣行了这事？}若瑟说：\Quote{我知道耶稣是一个公义、诚实、全然良善的人；我也知道你们因仇恨设法杀害了祂，因此我埋葬了祂。}于是司祭长大怒，捉住若瑟，把他关进监狱，对他说：\Quote{若不是明天是无酵节，我等明日也必像处死祂那样处死你；但你暂被拘留，到主日清晨，你必被处死。}他们这样说，便封了监狱，用各种锁链牢固锁住。\par

因此，当预备日结束，安息日清晨，犹太人往彼拉多那里去，对他说：\Quote{我主，那骗子曾说，三天后祂要复活。因此，恐怕祂的门徒夜间将祂偷去，用这样的欺骗迷惑百姓，请下令守护坟墓。}彼拉多于是给了他们五百名士兵，他们围在坟墓四周看守，并在墓石上加封了印。\BibleRef{玛 27:62}\par

主日黎明时，司祭长同犹太人召开了公议会，派人去提若瑟出狱，要处死他。但他们打开监狱，却找不到他。他们对此惊异不已——门关着，门闩完好，印封未破，若瑟竟消失了。\par

% structure: p0062 source=h2 target=subsection reason=relative-heading-rank; base=h2

\subsection{第十三章}

这时，看守坟墓的一个兵士上前来，在会堂里说：\Quote{你们要知道，耶稣已经复活了。}犹太人说：\Quote{怎么一回事？}他说：\Quote{先是发生地震；然后主的一位天使，身披闪电，从天降下，把石头从坟墓滚开，坐在上面。我们所有兵士因怕他，都变得像死人一样，不能逃跑，也不能说话。我们听见天使对那些来看坟墓的妇女说：\InnerQuote{你们不要害怕，我知道你们在寻找耶稣。他不在这里，已经复活了，正如他先前告诉你们的。你们俯身看看他身体安放的地方，但是去告诉他的门徒，他已从死者中复活了，让他们到加里肋亚去，因为在那里他们要找到他。}因此我首先告诉你们这件事。}\BibleRef{玛 28:1}\par

犹太人对兵士说：\Quote{那些到坟墓来的妇女是什么样的人？你们为什么没有抓住她们？}兵士说：\Quote{由于害怕和看见天使，我们既不能说话，也不能动弹。}犹太人说：\Quote{我指永生的以色列天主起誓，我们一句也不相信你们所说的话。}兵士说：\Quote{耶稣行了这样大的奇事，你们都不相信，你们会相信我们吗？你们说得对，天主是永生的；而你们所钉死的那一位确实是永生的。但我们听说你们曾把若瑟关在监牢里，后来你们开了门，却没有找到他。你们把若瑟带来，我们也就把耶稣带来。}犹太人说：\Quote{那个从监狱逃跑的若瑟，你们可以在他的故乡阿黎玛特雅找到他。}兵士说：\Quote{你们也到加里肋亚去，你们会找到耶稣，正如天使对妇女所说的那样。}\par

听了\Emph{这些话}，犹太人就害怕，对兵士说：\Quote{你们要当心，不要把这件事告诉任何人，否则人人都要相信耶稣了。}因此他们也给了他们许多钱。兵士说：\Quote{我们怕万一比拉多听说我们收了钱，他会杀死我们。}犹太人说：\Quote{收下吧；我们保证会在比拉多面前为你们辩护。你们只要说，你们当时睡着了，耶稣的门徒趁你们熟睡时，来把他从坟墓偷走了。}于是兵士拿了钱，就照他们所吩咐的说了。直到今日，这谎言仍在犹太人中间流传。\BibleRef{玛 28:11}\par

% structure: p0066 source=h2 target=subsection reason=relative-heading-rank; base=h2

\subsection{第十四章}

过了几天，有三人从\Place{加里肋亚}来到\Place{耶路撒冷}。其中一人是司祭，名叫\Person{丕乃哈斯}；第二人是肋未人，名叫\Person{阿盖}；第三人是兵士，名叫\Person{阿达斯}。他们来到大司祭们面前，对他们和民众说：\Quote{你们所钉死的耶稣，我们在\Place{加里肋亚}看见他同他的十一位门徒在\Place{橄榄山}上，教导他们说：\InnerQuote{你们往普天下去，宣讲福音；信而受洗的，必得救；不信的，必被判罪。}说了这话，他就被提升天了。}\BibleRef{谷 16:16}我们和另外五百多人都亲眼看见。\BibleRef{格前 15:6}\par

大司祭们和犹太人听了这些话，就对这三个人说：\Quote{你们要光荣以色列的天主，并要悔改你们所说的这些谎言。}他们回答说：\Quote{我指我们祖先亚巴郎、依撒格和雅各伯永生的天主起誓，我们没有说谎，而是告诉你们真理。}于是大司祭发言，有人从圣殿里拿出希伯来人的旧\Emph{盟约}，他使他们起誓，又给了他们钱，打发他们到别处去，免得他们在耶路撒冷宣扬基督的复活。\par

当众人听见这些事，群众就聚集到圣殿里，起了很大的骚动。因为许多人说：\Quote{耶稣已经从死者中复活了，正如我们所听见的，你们为什么钉死他？}\Person{亚纳斯}和\Person{盖法}说：\Quote{不要相信，犹太人哪，不要相信兵士说的话；也不要相信他们看见天使从天降下。因为我们给了兵士钱，叫他们不要把这些事告诉任何人；同样，耶稣的门徒也给了他们钱，叫他们说耶稣已经从死者中复活了。}\par

% structure: p0070 source=h2 target=subsection reason=relative-heading-rank; base=h2

\subsection{第十五章}

\Person{尼苛德摩}说：\Quote{耶路撒冷的居民啊，\Person{厄里亚}先知曾乘火马车升到高天，所以耶稣复活也并非难以置信；因为\Person{厄里亚}先知是耶稣的预像，为使你们听见耶稣复活而不致不信。因此我提议并劝告，我们应当派兵士到\Place{加里肋亚}去，到这些人作证说看见他和他的门徒的地方，好让他们四处寻找，找到他，这样我们可以因我们对他的恶行而向他求宽恕。}这个建议令他们满意；他们就选了兵士，打发他们往\Place{加里肋亚}去。他们确实没有找到耶稣，却找到了在\Place{阿黎玛特雅}的若瑟。\par

兵士回来后，大司祭们得知若瑟已被找到，就召集民众，说：\Quote{我们怎样做才能使若瑟到我们这里来？}他们商议后，就给他写了一封信，内容如下：\Quote{若瑟父亲，愿平安\Emph{归于}你和你的全家以及你的朋友们！我们知道我们得罪了天主，也得罪了你——他的仆人。因此，我们恳求你到你这些儿女这里来。因为我们很惊奇你是怎样从监狱里逃出来的，我们真实地说，我们曾对你怀有恶计。但是天主看到我们对你的谋划是不义的，就把你从我们手中救出。但请你到我们这里来，因为你是我们百姓的光荣。}\par

犹太人把这封信送到\Place{阿黎玛特雅}，由七个与若瑟为友的兵士带去。他们去了，找到了他；按照吩咐，恭敬地向他请安，把信交给他。他接过信，读了之后，就光荣天主，并拥抱兵士；然后摆上筵席，与他们整日整夜一同吃喝。\par

次日，他同他们一起出发前往耶路撒冷；百姓出来迎接他，拥抱了他。\Person{尼苛德摩}把他接进自己家中。又过了一天，\Person{亚纳斯}和\Person{盖法}，大司祭们，在圣殿里召见他，对他说：\Quote{你要光荣以色列的天主，把真相告诉我们。因为我们知道你埋葬了耶稣；为此我们抓住了你，把你关在监狱里。以后，当我们想把你带出来处死时，却找不到你，我们非常惊讶和害怕。再者，我们向天主祈祷，好能找到你，并问你。所以，请把真相告诉我们。}\par

\Person{若瑟}对他们说：\Quote{在预备日晚上，你们把我关在狱中时，我整夜祈祷，整个安息日也如此。到了半夜，我看见那监狱被四个天使抬了起来，他们抓住四角。耶稣像闪电一样进来，我因恐惧倒在地上。于是他拉着我的手，扶我起来，说：\InnerQuote{不要害怕，\Person{若瑟}。}接着，他拥抱我，亲吻我，并说：\InnerQuote{你转过身来，看看我是谁。}我转过身来，看着他说：\InnerQuote{我主，我不知道你是谁。}他说：\InnerQuote{我就是耶稣，你前天埋葬的那一位。}我对他说：\InnerQuote{请你让我看看坟墓，然后我就相信。}于是他拉着我的手，领我到那已打开的坟墓。看见殓布和汗巾，认出他来，我便说：\BibleQuote{奉主名而来的，当受赞美}；并朝拜了他。然后，他拉着我的手，由天使陪伴，带我回到我在\Place{阿黎玛特雅}的家中，对我说：\InnerQuote{你在这里坐四十天；因为我要到我的门徒那里去，好使他们能完整地宣讲我的复活。}}\par

% structure: p0076 source=h2 target=subsection reason=relative-heading-rank; base=h2

\subsection{第十六章。}

\Person{若瑟}这样说了之后，大司祭们向百姓呼喊说：\Quote{我们知道耶稣有父亲和母亲；我们怎能相信他就是基督？}有一个肋未人回答说：\InnerQuote{我知道耶稣的家族，他们是高尚的人，是天主的伟大仆人，并从犹太人中收取什一奉献。我也知道长老\Person{西默盎}，他在耶稣还是婴儿时接待了他，并对他说：\BibleQuote{主啊，现在可让你的仆人平安离去。}}\par

犹太人说：\Quote{现在我们去找那三个在\Place{橄榄山}上看见他的人，好能问问他们，更准确地了解真相。}他们找到了他们，把他们带到众人面前，让他们发誓说实话。他们说：\Quote{正如以色列的天主永在那样，我们看见耶稣活着在\Place{橄榄山}上，并升天去了。}\par

于是\Person{亚纳斯}和\Person{盖法}把三个人分别带到一旁，单独一个接一个地问他们。他们彼此一致，甚至三个人都给出了同样的说法。大司祭们回答说，说：\Quote{我们的圣经说，\BibleQuote{凭两三个见证人，一切陈述都应成立}。那么，\Person{若瑟}已经承认，他和\Person{尼苛德摩}一起料理了他的身体并埋葬了他，而且他确实复活了。}\par

\SourceRecord{Translated by Alexander Walker.；Ante-Nicene Fathers；Vol. 8.；Edited by Alexander Roberts, James Donaldson, and A. Cleveland Coxe.；Buffalo, NY: Christian Literature Publishing Co.,；1886.；<http://www.newadvent.org/fathers/08071b.htm>.}

% structure: p0001 source=h1 target=section reason=page-title

\section{\WorkTitle{尼苛德摩福音（第一部，拉丁文形式）}}

\Emph{比拉多行传}\par

我，\Person{厄内阿}，起初是希伯来人的保护者，并遵守法律；后来救主的恩宠和祂的伟大恩赐占据了我。我因圣经认识了基督耶稣，就来到祂面前，信奉了祂的信德，以便能堪当领受祂的圣洗圣事。首先，我寻找了有关我们的主耶稣基督在那个时代所写的记录，就是犹太人在本丢·比拉多时代所公布的，我们在希伯来文的作品中找到了它们，是在主耶稣基督的时代编写的；我在杰出的德奥多西皇帝执政第十七年，瓦伦提尼安第五次执政，第九小纪时，将它们译成了外邦人的语言。你们中凡阅读这本书并传抄者，请纪念我，为我祈祷，\Person{厄内阿}——天主最卑微的仆人，愿祂怜悯我，并宽恕我得罪祂的一切罪过。愿平安归于所有阅读这些的人，以及他们全家，直到永远！阿们。\par

当时，在罗马皇帝提庇留·凯撒在位第十九年，加里肋亚王黑落德之子黑落德统治第十九年，三月二十五日，即四月朔日前第八天，鲁菲诺和鲁贝利奥执政，第202届奥林匹克运动会第四年，犹太人的司祭若瑟和盖法治理时期：大司祭和其他犹太人所行的事，尼苛德摩在主受十字架和苦难之后记录下来，尼苛德摩本人用希伯来文记录下来。\par

% structure: p0005 source=h2 target=subsection reason=relative-heading-rank; base=h2

\subsection{第一章}

\Person{亚纳斯}和\Person{盖法}、\Person{苏玛斯}和\Person{达堂}、\Person{加玛里耳}、\Person{犹达斯}、\Person{肋未}、\Person{纳斐塔里}、\Person{亚历山大}和\Person{雅依洛}，以及其余的犹太人，来到比拉多面前，控告主耶稣基督许多事，说：\Quote{我们知道他是木匠若瑟的儿子，生于玛利亚；他说他是天主子，又是王。不仅如此，他还违犯安息日，并且想要废除我们祖传的法律。}比拉多说：\Quote{他所做的事是什么？他想要废除法律？}犹太人说：\Quote{我们有法律，不可在安息日医治任何人；但他用邪术在安息日治好瘸子、驼背、瞎子、瘫痪者、麻风病人和附魔的人。}比拉多对他们说：\Quote{凭藉什么邪术？}他们对他说：\Quote{他是巫师；借魔王贝耳则步驱魔，众魔都服从他。}比拉多对他们说：\Quote{驱魔不是靠邪神，而是靠医神斯克洛皮乌斯。}\par

犹太人说：\Quote{我们恳请陛下将他带到你的审判台前受审。}比拉多召集犹太人，对他们说：\Quote{我身为总督，怎能审判一个王呢？}他们对他说：\Quote{我们不说他是王，是他自己说他是。}比拉多召来一个差役，对他说：\Quote{用温和的态度把耶稣带进来。}差役出去，认出了祂，便朝拜了祂，并把手中拿的外衣铺在地上，说：\Quote{我主，请踏着这进来，因为总督叫您。}但犹太人看到差役所做的事，就向比拉多喊叫说：\Quote{你为什么不用传令官的声音叫他进来，而用差役？因为那差役看见他也朝拜了他，并把手中拿的外衣铺在他面前地上，对他说：\InnerQuote{我主，总督叫您。}}\par

比拉多召来差役，对他说：\Quote{你为什么这样做，并尊敬那称为基督的耶稣？}差役对他说：\Quote{当你派我去耶路撒冷找\Person{亚历山大}时，我看见他骑着一匹驴，希伯来儿童从树上折下树枝，铺在路上；其他人手拿树枝；另一些人将衣服铺在路上，呼喊说：\BibleQuote{求你拯救，因此，你是在至高之天的；\Emph{奉主名而来的是应当称颂的}！}}\par

犹太人大声喊叫，反对差役说：\Quote{希伯来儿童确实是用希伯来语喊的。你一个外邦人，怎么能知道？}差役对他们说：\Quote{我问了一个犹太人，说：\InnerQuote{他们用希伯来语喊的是什么？}他就给我解释了。}比拉多对他们说：\Quote{他们是怎么用希伯来语喊的？}犹太人说：\Quote{贺三纳于至高之天！}比拉多对他们说：\Quote{\Quote{贺三纳于至高之天}是什么意思？}他们对他说：\Quote{求你拯救我们，你是在至高之天的。}比拉多对他们说：\Quote{如果你们自己为儿童所喊的词语和话语作证，那么差役有什么罪？}他们就沉默了。总督对差役说：\Quote{出去，随你的意把他领进来。}差役出去，照先前一样做了，对耶稣说：\Quote{我主，请进，因为总督叫您。}\par

当时耶稣正进去，执旗者举着旗标，旗标的头像自动低下，朝拜了耶稣。犹太人看见旗标如何弯下并朝拜耶稣，就更加对执旗者喊叫。比拉多对犹太人说：\Quote{你们不惊讶旗标自动弯下并朝拜耶稣吗？}犹太人对比拉多说：\Quote{我们看见那扛旗标的人自己弯下并朝拜耶稣。}总督召来执旗者，对他们说：\Quote{你们为什么这样做？}他们对比拉多说：\Quote{我们是外邦人，是神庙的奴隶，我们怎么会朝拜他？因为我们拿着偶像时，它们自己弯下并朝拜了他。}\par

比拉多对会堂长和民间的长老说：你们挑选强壮有力的人来拿旗帜，让我们看看它们会不会自己鞠躬。犹太人的长老就选了十二个极强壮有力的人，叫他们拿旗帜，六人一面；他们站在总督的审判座前。比拉多对差役说：把耶稣带到总督府外面，然后随你的意思再带他进来。耶稣和差役就出了总督府。比拉多叫来先前拿旗帜的人，对他们说：我凭凯撒的健康发誓，如果耶稣进来时这些旗帜不自己鞠躬，我就砍你们的头。总督又吩咐耶稣第二次进来。差役照先前的样子，恳切求耶稣踩在他的袍子上走过去。耶稣就踩在上面，走了进去。耶稣进去时，旗帜立刻鞠躬朝拜了耶稣。\par

% structure: p0012 source=h2 target=subsection reason=relative-heading-rank; base=h2

\subsection{第二章}

比拉多见了，恐惧抓住他，立刻想要从审判座上起来。他正想着要起身离开的时候，他的妻子派人来对他说：\Quote{不要干涉那义人的事，因为我今夜为他受了许多苦。}比拉多就召来犹太人对他们说：你们知道我的妻子是敬拜天主的人，而且在犹太教方面更认同你们。犹太人对他说：的确如此，我们知道。比拉多对他们说：看，我的妻子派人来对我说：\Quote{不要干涉那义人的事，因为我今夜为他受了许多苦。}犹太人回答说：我们不是对你说过他是术士吗？看，他把梦象差遣给了你的妻子。\par

比拉多叫来耶稣，对他说：这些人作证反对你的事，你什么都不回答吗？耶稣回答说：他们若没有权柄，就不会说话。各人有权柄用自己的口说善说恶；让他们自己瞧着办吧。\par

犹太人的长老回答耶稣说：我们要瞧着办什么？第一，你是生于淫乱的；第二，你出生在白冷时发生了屠杀婴儿的事；第三，你的父亲若瑟和你的母亲玛利亚逃往埃及，因为他们对百姓没有信任。\par

旁观的几个善心的犹太人说：我们说他不是生于淫乱的；我们知道玛利亚已经许配给若瑟，他不是生于淫乱的。比拉多对说他生于淫乱的犹太人说：你们这话不真，因为订婚已经发生了，像你们本族这些人所说的。亚纳斯和盖法对比拉多说：我们和全体群众都说他是生于淫乱的，而且他是术士；但这些人都是归化者\OriginalTerm{归化者}{proselytes}，是他的门徒。比拉多叫来亚纳斯和盖法，对他们说：什么是归化者？他们对他说：他们生为外邦人的儿子，后来成了犹太人。于是那些作证耶稣不是生于淫乱的人——拉匝禄、阿斯泰里、安多尼、雅各伯、安纳斯、阿匝辣、撒慕尔、依撒格、斐尼、克里斯普、阿格黎帕、犹大——回答说：我们不是生为归化者，而是犹太人的儿子，我们说实话；因为玛利亚订婚时我们在场。\par

比拉多叫来那十二个证明耶稣不是生于淫乱的人，对他们说：我凭凯撒的健康命令你们，告诉我耶稣不是生于淫乱是不是真的。他们对比拉多说：我们有法律不可发誓，因为那是罪；但让他们凭凯撒的健康发誓说我们说的不对，我们就该当死罪。比拉多就对亚纳斯和盖法说：这些人所证明的，你们什么都不回答吗？亚纳斯和盖法对比拉多说：那十二个人被相信说他不生于淫乱；我们——全体百姓——喊叫说他是生于淫乱的，是术士，并且自称是天主子、君王，我们却不被相信。\par

比拉多吩咐全体群众出去，只留下那十二个说他不生于淫乱的人，并吩咐把耶稣与他们分开。比拉多对他们说：犹太人为什么要处死耶稣？他们对他说：他们生气因为他在安息日治病。比拉多说：因一件善事他们就要处死他吗？他们对他说：是的，大人。\par

% structure: p0019 source=h2 target=subsection reason=relative-heading-rank; base=h2

\subsection{第三章}

比拉多满怀怒气，走到总督府外面对他们说：我以太阳作证，我在这人身上找不出一点罪来。犹太人回答总督说：他若不是作恶的，我们就不会把他交给你。比拉多对他们说：你们带他去，按你们的法律审判他。犹太人回答说：我们不可处死人。比拉多对他们说：天主吩咐你们不可处死人吗？难道他就吩咐要我杀人吗？\par

比拉多又进了总督府，私下叫耶稣来对他说：你是犹太人的王吗？耶稣回答比拉多说：你这话是自己说的，还是别人对你论我说的？比拉多回答说：我难道是犹太人吗？你的民族和司祭长把你交给了我。你做了什么？耶稣回答说：我的国不属于这世界。若是我的国属于这世界，我的臣仆必要努力使我不被交给犹太人；但现在我的国不从此地而来。比拉多对他说：那么你是王吗？耶稣对他说：你说我是王。我为此而生，为此而来，要为真理作证；凡属于真理的，都听我的声音。比拉多对他说：什么是真理？耶稣说：真理是从天上来的。比拉多说：地上难道没有真理吗？耶稣对比拉多说：请注意，说真理的人如何受地上有权柄者的审判。\par

% structure: p0022 source=h2 target=subsection reason=relative-heading-rank; base=h2

\subsection{第四章}

比拉多于是将耶稣留在总督府内，出来见犹太人，对他们说：\Quote{我在这人身上查不出丝毫过错。}犹太人对他说：\Quote{他说过：\InnerQuote{我能拆毁这座圣殿，并在三天内把它重建起来。}}比拉多对他们说：\Quote{什么圣殿？}犹太人对他说：\Quote{\Emph{那座圣殿}，是撒罗满用了四十六年建造的；而他说\Emph{他能}在三天内拆毁并重建它。}比拉多对他们说：\Quote{对这人的血，我是无罪的；你们\Emph{自己}看着办吧。}犹太人对他说：\Quote{他的血\Emph{归}在我们和我们子孙身上。}\par

比拉多又召集了长老、司祭和肋未人，私下对他们说：\Quote{不要这样做；因为你们控告他的一切，我查不出他该死之罪，连关于治愈和干犯安息日的事也没有。}司祭、肋未人和长老说：\Quote{请问，如果有人亵渎凯撒，他该不该死？}比拉多说：\Quote{他该死。}犹太人回答他：\Quote{那么，那亵渎天主的，岂不更该死吗！}\par

总督命犹太人走出总督府；然后叫来耶稣，对他说：\Quote{我该对你怎么办？}耶稣对比拉多说：\Quote{\Emph{你}已经得到了。}比拉多说：\Quote{怎么得到了？}耶稣说：\Quote{梅瑟和先知们已预告了我的死亡和复活。}犹太人听见了，对比拉多说：\Quote{你为什么还要听这种亵渎的话？}比拉多说：\Quote{如果这话是亵渎，你们就把他带走，带到你们的会堂，按你们的法律审判他。}犹太人对比拉多说：\Quote{我们的法律规定：一个人若得罪了人，该受四十减一的鞭刑；但若亵渎了天主，该被石头砸死。}\par

比拉多对他们说：\Quote{那么就按你们的法律审判他吧。}犹太人对比拉多说：\Quote{我们要求把他钉死在十字架上。}比拉多对他们说：\Quote{他不该被钉十字架。}\par

总督环视周围站着的犹太民众，看见许多犹太人哭泣，便说：\Quote{全体民众都不愿他死。}长老们对比拉多说：\Quote{我们——全体民众——正是为此而来，要他死。}比拉多对犹太人说：\Quote{他做了什么该死的事？}他们说：\Quote{因为他自称是天主子，又是君王。}\par

% structure: p0028 source=h2 target=subsection reason=relative-heading-rank; base=h2

\subsection{第五章}

有一个叫尼苛德摩的犹太人站在总督面前，说：\Quote{求求你，请允许我说几句话。}比拉多对他说：\Quote{说吧。}尼苛德摩说：\Quote{我在会堂中对长老、司祭、肋未人和全体犹太民众说过：你们与这人\Emph{有何关系}？这人行了许多奇事和征兆，是没有人做过或能做到的。放了他，不要对他图谋不轨：他所行的征兆若是出于天主，就会存立；若是出于人，就会归于无有。因为梅瑟奉天主差遣进入埃及，也行了天主吩咐他在埃及王法郎面前行的许多征兆。当时有术士雅乃斯和杨布勒在那里施行法术，他们也做了梅瑟所行的征兆，但并非全部；埃及人竟把他们当作神明，就是雅乃斯和杨布勒。由于他们所行的征兆不是出于天主，他们便灭亡了，连同那些信从他们的人。现在让这人走吧，因为他并不该死。}\par

犹太人对尼苛德摩说：\Quote{你成了他的门徒，为他辩护。}尼苛德摩对他们说：\Quote{难道总督也成了他的门徒，为他辩护吗？凯撒不是派他担任这个职位吗？}犹太人向尼苛德摩发怒，咬牙切齿。比拉多对他们说：\Quote{你们\Emph{正听见}真理，为什么还向他咬牙切齿？}犹太人对尼苛德摩说：\Quote{愿你承受他的真理，和他同分！}尼苛德摩说：\Quote{阿们，阿们，阿们；愿我承受\Emph{它}，正如你们所说的！}\par

% structure: p0031 source=h2 target=subsection reason=relative-heading-rank; base=h2

\subsection{第六章}

又有一个犹太人站起来，请求总督允许他说句话。总督说：\Quote{你说吧。}他说：\Quote{我卧病在床三十八年，痛苦不堪。耶稣来时，许多附魔的人和被各种疾病缠身的人都被他治愈了。有几个年轻人怜悯我，抬着我的床，把我放在他面前。耶稣看见我，就怜悯我，对我说话：\InnerQuote{拿起你的床，走吧。}我立刻痊愈了，拿起床走了。}犹太人对比拉多说：\Quote{问他是在哪一天被治愈的。}他说：\Quote{安息日。}犹太人说：\Quote{我们不是告诉过你，他在安息日治病，驱逐魔鬼吗？}\par

又有一个犹太人站起来说：\Quote{我生来瞎眼；我听见声音，却看不见人。耶稣路过时，我大声呼喊：\InnerQuote{达味之子，可怜我吧！}他怜悯了我，把手放在我的眼睛上，我立刻就看见了。}又有一个犹太人站起来说：\Quote{我驼背，他一句话就使我直了。}另一个说：\Quote{我患了麻风病，他一句话就治愈了我。}\par

% structure: p0034 source=h2 target=subsection reason=relative-heading-rank; base=h2

\subsection{第七章}

又有一个名叫韦罗尼加的女人，远远地向总督喊道：\Quote{我患血漏十二年；我一摸他衣服的穗边，血漏立刻就止住了。}犹太人说：\Quote{我们有法律规定，女人不可作证。}\par

% structure: p0036 source=h2 target=subsection reason=relative-heading-rank; base=h2

\subsection{第八章}

又有许多其他男女大声说：\Quote{那人是先知，魔鬼都服从他。}比拉多对那些说魔鬼服从他的人说：\Quote{你们的师父为什么不服从他呢？}他们对比拉多说：\Quote{我们不知道。}另有人对比拉多说：\Quote{他使死了四天的拉匝禄从坟墓里复活了。}总督听了，战栗地对全体犹太民众说：\Quote{你们为什么要流无辜人的血？}\par

% structure: p0038 source=h2 target=subsection reason=relative-heading-rank; base=h2

\subsection{第九章}

比拉多叫了尼苛德摩和那十二个说祂不是由奸淫所生的人来，对他们说：\Quote{百姓起了骚动，我该怎么办？}他们对他说：\Quote{我们不知道；让他们自己看着办吧。}比拉多又召集全体犹太群众，说：\Quote{你们知道，在无酵节那天，我有惯例释放一个囚犯给你们。我监狱里有一个出名的囚犯，是杀人犯，名叫巴辣巴；还有耶稣，称为基督，我在祂身上找不出死因。你们愿意我释放哪一个给你们？}他们全都喊叫说：\Quote{释放巴辣巴给我们！}比拉多对他们说：\Quote{那么，我怎样处置那称为基督的耶稣呢？}他们都说：\Quote{让祂被钉十字架。}犹太人又说：\Quote{你如果释放这个人，就不是凯撒的朋友，因为祂自称是天主子，又是君王；除非你愿意这个人作王，而不是凯撒。}\par

比拉多满腔愤怒，对他们说：\Quote{你们的民族总是叛逆的，你们总是反对那些为你们好的人。}犹太人回答说：\Quote{谁是拥护我们的？}比拉多对他们说：\Quote{你们的天主，祂把你们从埃及的残酷奴役中救出，领你们出埃及，经海如走干地，在旷野用玛纳和鹌鹑喂养你们，从岩石中给你们流出水来，赐给你们喝，并赐给你们法律；而你们在这一切事上激怒了你们的天主，为自己寻求一个神，一个铸成的牛犊。你们激怒了你们的天主，祂想要杀死你们；梅瑟为你们求情，使你们不致灭亡。现在你们说朕憎恨君王。}\par

比拉多从审判台上起来，想要出去。犹太人喊叫说：\Quote{我们知道凯撒是王，耶稣不是。因为贤士们也向祂献了礼物，如同向王一样；黑落德从贤士们那里听说一位王诞生，就想要杀祂。这事被知道后，祂的父亲若瑟就带着祂和祂的母亲逃往埃及；黑落德听说后，就屠杀了白冷出生的犹太婴孩。}\par

比拉多听了这些话，害怕起来。在喊叫的群众中安静下来后，比拉多说：\Quote{那么，这就是黑落德所寻找的那一位吗？}他们对他说：\Quote{正是祂。}比拉多拿水在群众面前洗手，说：\Quote{对这义人的血，我是无罪的；你们自己看着办吧。}犹太人又喊叫说：\Quote{祂的血归在我们和我们子孙身上。}\par

于是比拉多下令松开帷幔，对耶稣说：\Quote{你的子民控告你是王；因此朕首先按照皇帝的律法判决你受鞭打，然后被钉在十字架上。}\par

% structure: p0044 source=h2 target=subsection reason=relative-heading-rank; base=h2

\subsection{第十章}

耶稣受鞭打后，比拉多把祂交给犹太人钉十字架，同钉的还有两个强盗；一个名叫\OriginalTerm{狄斯玛}{Dismas}，另一个名叫\OriginalTerm{革斯达}{Gestas}。他们到了那地方，就剥去祂的衣服，给祂围上一条麻布，把一顶荆棘冠冕戴在祂头上。同样他们也把两个强盗同祂一起钉了，狄斯玛在右边，革斯达在左边。耶稣说：\BibleQuote{父啊，宽赦他们吧，因为他们不知道他们做的是什么。}士兵们分了祂的衣服。百姓站着观望；他们的司祭长和判官嘲笑祂，彼此说：\Quote{祂救了别人，如今让祂救自己吧；如果祂是天主子，就让祂从十字架上下来吧！}士兵们也嘲笑祂，俯伏在祂面前，拿醋和苦胆给祂喝，说：\Quote{如果你是犹太人的王，就救自己吧。}\par

比拉多在判决后，下令用希伯来、希腊和拉丁文字写一个牌子，按照犹太人的话：\Quote{这是犹太人的王。}\par

被钉的一个强盗，名叫革斯达，对祂说：\Quote{如果你是基督，救你自己和我们吧。}狄斯玛回答说，斥责他说：\Quote{你既然受同样的刑罚，连天主都不怕吗？我们受刑是公平的，因为我们所忍受的，是我们应得的；但祂没有做过一件坏事。}他不断对耶稣说：\Quote{主，当祢来临时，请记得我。}耶稣对他说：\BibleQuote{我实在告诉你，今天你就要同我在乐园里。}\par

% structure: p0048 source=h2 target=subsection reason=relative-heading-rank; base=h2

\subsection{第十一章}

约第六时辰，黑暗笼罩了整个大地；太阳昏暗，圣所的帐幔从中间裂开。耶稣大声呼喊说：\BibleQuote{父啊，我把我的灵魂交在祢手里。}说完这话，就断了气。百夫长看见所发生的事，光荣天主说：\BibleQuote{这真是个义人。}所有在场观看的群众，看见所发生的事，都捶着胸回去了。\par

百夫长向总督报告了所发生的事。总督和他的妻子听见，非常悲伤，那天不吃不喝。比拉多召集犹太人对他们说：\Quote{你们看见所发生的事了吗？}他们对总督说：\Quote{发生了日蚀，照常如此。}\par

祂的相识们以及从加里肋亚跟随祂的妇女们远远站着观看这些事。有一个人，名叫若瑟，是个议员，为人善良正义，没有赞同他们的计谋和行为，是从犹太的阿黎玛特雅城来的，他也等待天主的国，去见比拉多，请求领取耶稣的遗体。他把祂从十字架上卸下来，用干净的殓布裹好，安放在他自己的新坟墓里，这坟墓是从未安放过人的。\par

% structure: p0052 source=h2 target=subsection reason=relative-heading-rank; base=h2

\subsection{第十二章}

犹太人听说\Person{若瑟}请求了\Person{耶稣}的身体，就四处搜寻他；还有那十二个说祂不是由淫乱而生的人，以及\Person{尼苛德摩}，和许多其他曾在\Person{比拉多}面前为祂作证、宣告祂善行的人。他们都躲藏起来，只有\Person{尼苛德摩}独自出现在他们面前，因为他是犹太人的首领；他对他们说：你们怎么进了会堂？犹太人对他说：你既然赞同他，又怎么进了会堂？愿他的分与你在来世同在！\Person{尼苛德摩}说：阿们，阿们，阿们。同样，\Person{若瑟}也出来对他们说：你们为什么向我发怒，因为我求了\Person{耶稣}的身体？看，我把他安放在我的新坟墓里，用洁净的细麻布包裹，又滚了一块石头堵住洞口。你们对待一个义人并不好，因为你们没有记住你们如何把他钉在十字架上，并用长矛刺透了他。于是犹太人抓住\Person{若瑟}，因安息日的缘故，下令将他囚禁；他们对他说：你要知道，因时辰所迫，我们不能对你做什么，因为安息日快要到了。但你要明白，你连埋葬都不配，我们要把你的肉给空中的飞鸟和地上的野兽。\Person{若瑟}对他们说：那是骄傲的\Person{哥肋雅}说的话，他辱骂了生活的天主，反对圣\Person{达味}。天主说过：\BibleQuote{伸冤在我，我必报应，主这样说。}\Person{比拉多}心中不安，取水在太阳面前洗手，说：\BibleQuote{我流这义人的血，罪不在我；你们自己看着办吧。}你们却回答\Person{比拉多}说：\BibleQuote{他的血归在我们和我们子孙身上。}现在我怕天主的忿怒总有一天会临到你们和你们的子孙身上，正如你们所说的。犹太人听了这话，心中忿恨；他们抓住\Person{若瑟}，把他关在一个没有窗户的房子里，在门口设了守卫，并封了\Person{若瑟}被关的门。\par

到了安息日早晨，他们与司祭和肋未人商议，要在安息日后全体聚集。黎明醒来，会堂里的全体群众商议要用什么死法处死他。当集会坐定，他们命令以极大的侮辱把他带来；但打开门，却找不到他。于是众人都惊慌，极其惊讶，因为他们发现封条依然完好，而且\Person{盖法}拿着钥匙。他们再也不敢伸手去碰那些在\Person{比拉多}面前为\Person{耶稣}辩护的人。\par

% structure: p0055 source=h2 target=subsection reason=relative-heading-rank; base=h2

\subsection{第十三章}

他们正坐在会堂里，互相指责关于\Person{若瑟}的事，这时有几个守卫来了，就是他们曾从\Person{比拉多}那里请求来看守\Person{耶稣}坟墓的人，以免祂的门徒来偷走祂。他们报告说，对会堂的领袖、司祭和肋未人讲述了所发生的事：怎样发生了一场大地震，我们看见一位主的天使从天降下，把石头从墓门滚开，坐在上面；他的容貌像闪电，衣服像雪。我们因害怕，就像死人一样。我们听见天使对来到坟墓的妇女说话，说：\BibleQuote{你们不要害怕！我知道你们寻找被钉死的耶稣；祂不在这里，祂已复活了，正如祂所说：你们来看安放主的地方。快去告诉祂的门徒，祂已从死者中复活了，并在你们之先往加里肋亚去，正如祂对你们所说的。}\par

犹太人说：他对哪些妇女说话？士兵说：我们不知道那些妇女是谁。犹太人说：在什么时辰？守卫说：半夜。犹太人说：你们为什么没有抓住她们？守卫说：我们因为害怕天使，像死人一样，当时都不指望再见到白昼；又怎么能抓住她们呢？犹太人说：\Emph{正如}主天主永在，我们不信你们。守卫对犹太人说：你们在那人身上看见了如此大的征兆，却仍不信；那么你们怎能相信我们，说主活着呢？因为你们确实发誓说主\Person{耶稣基督}活着。守卫又对犹太人说：我们听说你们把求\Person{耶稣}身体的\Person{若瑟}关在监牢里，并用你们的戒指封了门；打开时却找不到他。那么，把\Person{若瑟}还给我们，我们就给你们\Person{耶稣基督}。犹太人说：\Person{若瑟}回自己的城市\Place{阿黎玛特雅}去了。守卫对犹太人说：而\Person{耶稣}，我们听天使说，在\Place{加里肋亚}。\par

犹太人听见这些话，极其害怕，说：只怕有一天这话传出去，众人都会信\Person{耶稣}。犹太人彼此商议，拿出相当数量的银币给士兵，说：你们要说，我们睡觉的时候，祂的门徒来把祂偷走了。如果这事被总督听见，我们会说服他，使你们平安无事。士兵收了\Emph{银钱}，就照犹太人所嘱咐的说；这话就在众人中间传开了。\par

% structure: p0059 source=h2 target=subsection reason=relative-heading-rank; base=h2

\subsection{第十四章}

有一个司祭名叫\Person{菲内斯}，一个教师名叫\Person{阿达斯}，一个肋未人名叫\Person{厄吉亚}，从\Place{加里肋亚}来到\Place{耶路撒冷}，向会堂的领袖、司祭和肋未人报告，他们看见\Person{耶稣}坐在\Place{橄榄山上}，门徒与祂同在；那山名叫\Place{曼布勒}或\Place{玛肋客}。\Person{耶稣}对门徒说：\BibleQuote{你们往普天下去，向一切受造物宣讲天国的福音。信而受洗的必要得救；不信的必要被定罪。这些征兆将伴随着信的人：因我的名驱逐魔鬼，说新语言，手拿毒蛇，若喝了什么致命的毒物，也决不受害；按手在病人身上，必得好转。}\Person{耶稣}正这样对门徒说话时，我们看见祂被提升到天堂。\par

司祭、肋未人和长老对他们说：\Quote{请你们将光荣归于以色列的天主，并向祂宣认，你们所讲述的这些事是否都是你们亲耳听见、亲眼看见的。}那些作报告的人说：\Quote{我们祖先的上主天主——\Person{亚巴郎}的天主、\Person{依撒格}的天主、\Person{雅各伯}的天主是永生的，我们确实听见并看见了。}犹太人对他们说：\Quote{你们来就是为了告诉我们这些吗？还是来向天主祈祷？}他们说：\Quote{我们是来向天主祈祷的。}长老、司祭长和肋未人对他们说：\Quote{如果你们是来向天主祈祷，为什么在全百姓面前抱怨那个愚蠢的故事？}司祭\Person{丕乃哈斯}、经师\Person{阿达斯}和肋未人\Person{厄基亚}对会堂首领、司祭和肋未人说：\Quote{如果我们所说的这些话——我们所看见和听见的——是罪，那么请看，我们就在你们面前，请按你们眼中看为好的对待我们。}他们拿起法律书，命令他们此后不可将这话告诉任何人。然后他们给他们吃喝，并将他们送出城外，还给了他们银子和碎块，并派三个人护送他们直到\Place{加里肋亚}。\par

当那些人上到\Place{加里肋亚}后，犹太人彼此商议；会堂的首领关起门来，心中充满极大的愤怒，说：\Quote{在以色列发生的是什么征兆？}\Person{亚纳斯}和\Person{盖法}说：\Quote{你们为什么心里忧愁？难道我们要相信那些士兵，说上主的一位天使从天上下来，滚开了墓门的石头吗？\Emph{不}！而是祂的门徒给了看守坟墓的人许多金子，把\Person{耶稣}带走，并教他们这样说：\InnerQuote{你们要说：上主的一位天使从天上下来，滚开了墓门的石头。}难道你们不知道，犹太人若相信外邦人的一句话就是违法的吗？况且这些人从我们这里得了足够的金子，就照我们教导的说了。}\par

% structure: p0063 source=h2 target=subsection reason=relative-heading-rank; base=h2

\subsection{第十五章}

\Person{尼苛德摩}站起来，站在议会的中间说：\Quote{你们说得对。那些从\Place{加里肋亚}下来的人，难道不是敬畏天主的人、爱好和平、憎恶谎言的人吗？他们起誓说：\InnerQuote{我们看见\Person{耶稣}坐在\Place{玛默勒}山上，和祂的门徒在一起，祂在我们耳中教导他们。}并且看见祂被接到天上去了。没有人问他们：祂是怎样被接到天上去的？正如圣经所教导我们的，圣\Person{厄里亚}也被接到天上去了；\Person{厄里叟}大声呼喊，\Person{厄里亚}把他的羊皮披在\Person{厄里叟}身上；\Person{厄里叟}又用那件羊皮披在\Place{约但河}上，就过去了，来到\Place{耶里哥}。先知们的儿子们遇见他，对\Person{厄里叟}说：\InnerQuote{你的师父\Person{厄里亚}在哪里？}他说：\InnerQuote{祂已被接到天上去了。}他们对\Person{厄里叟}说：\InnerQuote{难道有灵把他攫去，丢在一座山上吗？不如让我们带着我们的仆人去找他。}他们说服了\Person{厄里叟}，他就和他们同去。他们找了三天三夜，没有找到他，因为他已被接去了。现在，诸位，请听我说，让我们派人到全以色列去查看，恐怕\Person{耶稣}也被接到某处，丢在一座山上。}这话使众人都喜悦。他们就派人到以色列的各山去寻找\Person{耶稣}，却没有找到祂；但找到了\Place{阿黎玛特雅}的\Person{若瑟}，没有人敢去抓他。\par

他们向长老、司祭和肋未人报告说：\Quote{我们走遍了以色列的各山，没有找到\Person{耶稣}；但在\Place{阿黎玛特雅}找到了\Person{若瑟}。}他们听到\Person{若瑟}的事，就欢喜，将光荣归于以色列的天主。会堂的首领、司祭和肋未人商议怎样派人去见\Person{若瑟}，就取了纸，写信给\Person{若瑟}说：\par

\Quote{愿你和你所有的一切平安！我们知道我们得罪了天主，也得罪了你；你曾向以色列的天主祈祷，祂就把你从我们手中救了出来。现在请屈尊来到你的父老和你的儿女这里，因为我们极其忧伤。我们都寻找过你——我们开了门，却没有找到你。我们知道我们曾向你策划了恶计，但上主使我们的计划在你身上落空了。你配得全体百姓的尊敬，\Person{若瑟}父啊。}\par

他们从全以色列中选了七位\Person{若瑟}的朋友，\Person{若瑟}也知道他们是友好的；会堂的首领、司祭和肋未人对他们说：\Quote{看，如果他接了信并读了，他一定会跟你们到我们这里来；但如果他不读信，你们就知道他对我们心怀恶意，那么就向他请安，然后回到我们这里。}他们祝福了他们，就打发他们去了。他们来到\Place{阿黎玛特雅}\Person{若瑟}那里，俯伏在地朝拜他，说：\Quote{愿你和你所有的一切平安！}\Person{若瑟}说：\Quote{愿你们和全以色列百姓平安！}他们就把信卷交给他。\Person{若瑟}接了读了，卷起信，赞美天主说：\Quote{愿上主天主\Emph{受}赞美，祂从流无辜人之血中救出了以色列；愿天主\Emph{受}赞美，祂派遣了祂的天使，用祂的翅膀庇护了我。}他亲吻了他们，为他们摆设了筵席；他们便吃喝，并在那里过夜。\par

早晨他们起来，\Person{若瑟}备好驴，与他们一同上路，来到圣城\Place{耶路撒冷}。全体百姓都出来迎接他们，欢呼说：\Quote{愿平安\Emph{临在}你来临之时，\Person{若瑟}父啊！}他回答他们说：\Quote{愿上主的平安\Emph{临在}全体百姓！}众人都亲吻了他。他们同\Person{若瑟}一起祈祷，看见他都感到害怕。\Person{尼苛德摩}带他到自己家里，设了大宴，请了\Person{亚纳斯}、\Person{盖法}、长老、司祭长和肋未人到家中。他们与\Person{若瑟}一同欢乐、吃喝，赞美了天主，然后各回各家。\Person{若瑟}留在\Person{尼苛德摩}家里。\par

次日，即预备日，司祭、会堂首领和肋未人清早起来，来到\Person{尼苛德摩}家中。\Person{尼苛德摩}迎接他们，对他们说：\Quote{愿你们平安！}他们对他说：\Quote{愿你和\Person{若瑟}平安，愿你和你的家及\Person{若瑟}的家平安！}\Person{尼苛德摩}便领他们进入他家。公议会坐定，\Person{若瑟}坐在\Person{亚纳斯}与\Person{盖法}之间，没有人敢说话。\Person{若瑟}对他们说：\Quote{你们为什么召我来？}他们向\Person{尼苛德摩}使眼色，要他向\Person{若瑟}说话。\Person{尼苛德摩}开口说：\Quote{\Person{若瑟}父，你知道可敬的经师、司祭和肋未人想听你的一句话。}\Person{若瑟}说：\Quote{请问吧。}\Person{亚纳斯}和\Person{盖法}拿起\WorkTitle{法律书}，恳求\Person{若瑟}说：\Quote{请你光荣以色列的天主，并向祂宣认，不要向我们隐瞒任何话。}他们又对他说：\Quote{我们很忧伤，因为你求了\Person{耶稣}的尸体，用洁白的殓布裹好，安放在坟墓里。因此我们将你关在一个没有窗户的屋里，锁上门并加封；在每周的第一天我们开门，却找不到你。因此我们非常忧伤，所有天主的子民都惊骇不已。所以派人请你来，现在请告诉我们发生了什么事。}\par

于是\Person{若瑟}说：\Quote{在预备日约第十时辰，你们把我关起来，我整个安息日都在里面。当半夜来临，我站着祈祷时，你们关我的那屋子四角被提起，我眼前有电光闪耀。我战栗地跌倒在地。然后有人从我倒下的地方扶起我，从头顶到脚给我浇了大量的水，用奇妙的香膏抹我的鼻孔，又用水擦我的脸，如同给我洗脸，并亲吻我，对我说：\InnerQuote{\Person{若瑟}，不要怕；睁开你的眼睛，看是谁对你说话。}我一看，看见\Person{耶稣}，就惊恐，以为是幻影。我以祈祷和诫命对他说话，他也对我说话。我对他说：\InnerQuote{你是\Person{厄里亚}辣彼吗？}他对我说：\InnerQuote{我不是\Person{厄里亚}。}我说：\InnerQuote{我主，你是谁？}他对我说：\InnerQuote{我是\Person{耶稣}，你从\Person{比拉多}那里求了祂的身体，用洁白的殓布裹好，把汗巾放在我脸上，把我安放在你的新坟墓里，又滚了石头在墓门口。}于是我对那对我说话的说：\InnerQuote{主，请指示我，我把你放在哪里。}他领我，指示了我安放他的地方，以及我给他裹的殓布，和我裹在他脸上的汗巾；我便知道那是\Person{耶稣}。他用手拉住我，把我放在我屋子中间，虽然大门关着，又把我放在床上，对我说：\InnerQuote{愿你平安！}他亲吻我，对我说：\InnerQuote{四十天不要出你的屋子；因为，看，我要往加里肋亚到我弟兄们那里去。}}\par

% structure: p0071 source=h2 target=subsection reason=relative-heading-rank; base=h2

\subsection{第十六章。}

会堂首领、司祭和肋未人听了\Person{若瑟}这些话，如同死人一般，倒在地上，禁食到第九时辰。\Person{若瑟}和\Person{尼苛德摩}恳求他们说：\Quote{起来，站直，吃点食物，安慰你们的灵魂，因为明天是主的安息日。}他们便起来，恳求主，然后吃喝，各回自己的家。\par

安息日那天，经师和夫子们坐着互相询问，说：\Quote{这临到我们身上的愤怒是什么？因为我们认识他的父母。}\Person{经师肋未}说：\Quote{我知道他的父母敬畏天主，从不离开祈祷，且每年三次交纳什一。耶稣出生时，他的父母带他到这个地方，向天主献上祭品和全燔祭。确实，伟大的经师西默盎把他抱在怀里，说：\InnerQuote{主啊，现在照你的话，让你的仆人平安离去；因为我亲眼看见了你的救恩，就是你在万民面前所预备的，为启示万邦的光，为你子民以色列的荣耀。}他又祝福他的母亲\Person{玛利亚}，说：\InnerQuote{关于这孩子，我要向你宣布一件事。}\Person{玛利亚}说：\InnerQuote{我主，请说。}\Person{西默盎}说：\InnerQuote{好。}他又说：\InnerQuote{看，这孩子被立为使以色列中许多人跌倒和兴起，并成为反对的记号；至于你，要有一把剑刺透你自己的灵魂，为叫许多人心中的思念显露出来。}}\par

犹太人对\Person{肋未}说：\Quote{你怎么知道这些事？}\Person{肋未}说：\Quote{你们不知道我从他那里学了法律吗？}公议会的人说：\Quote{我们想见你的父亲。}他们便寻找他的父亲，得到了消息；因为他说：\Quote{你们为什么不信我的儿子？有福的义人西默盎教了他法律。}公议会对辣彼\Person{肋未}说：\Quote{你所说的话是真的。}司祭长、会堂首领和肋未人彼此说：\Quote{来，我们打发人到加里肋亚去，到那三个曾来这里报告他的教导和他被接升天的人那里，让他们告诉我们，他们怎样看见他被接到天上。}这话众人都喜悦。于是他们派了三人往加里肋亚去；他们说：\Quote{去对\Person{辣彼阿达斯}、\Person{辣彼菲内斯}和\Person{辣彼厄吉亚}说：\InnerQuote{愿你们和你们的家人平安！}公议会关于耶稣进行了许多调查；因此我们奉命召你们到圣地耶路撒冷来。}\par

那些人到了加里肋亚，看见他们坐着默想法律。他们以平安问候他们。他们说：\Quote{你们为什么来？}使者说：\Quote{公议会召你们到圣城耶路撒冷去。}那些人听说公议会寻找他们，便向天主祈祷，与使者一同躺卧，与他们一同吃喝。早晨起来，他们平安地去了耶路撒冷。\par

次日公议会坐堂，询问他们说：\Quote{你们确实看见\Person{耶稣}坐在\Place{曼勃勒山}上教导门徒，并被接到天上吗？}\par

首先，教师阿达斯说：\Quote{我确实看见祂坐在玛默勒山上教训祂的门徒；有一片发光的云遮蔽了祂和祂的门徒，祂就升到天上去了；祂的门徒就俯伏在地祈祷。}他们又召来司祭非乃斯，问他说：\Quote{你看见耶稣被接升天是怎样呢？}他说的和前者相同。他们又召了第三位，拉比厄吉亚斯，问他，他说的也和第一、第二位相同。那时在公议会中的人说：\Quote{梅瑟法律规定，凭两个或三个人的口供，一切事都应成立。}教师之一，经师中的一位阿布德姆说：\Quote{法律上记载：哈诺客与天主同行，就被接去了，因为天主将他取去。}教师雅依洛说：\Quote{我们也听见过圣梅瑟的死，却未曾看见\Emph{它}；因为主的法律上记载：梅瑟按天主的话死了，直到今日没有人知道他的埋葬。}拉比肋未说：\Quote{西默盎长老所说的\InnerQuote{祂被立定，使以色列中许多人跌倒和兴起，并成为一个反对的记号}，这是什么意思呢？}拉比依撒格说：\Quote{法律上记载：\InnerQuote{看，我派遣我的使者，在你面前先行，保护你行一切善道，因为我已使祂的新名来到。}}\par

于是亚纳斯和盖法说：\Quote{你们说得对，这些事都记载在梅瑟法律上：没有人看见哈诺客的死，也没有人指出圣梅瑟的埋葬。而耶稣在比拉多面前受审，我们看见祂受鞭打，脸上被人吐唾沫；兵士们给祂戴上荆棘冠冕，祂从比拉多那里接受了判决；然后祂被钉在十字架上，他们给祂喝苦胆和醋，有两个强盗和祂一同被钉；士兵隆基诺用长枪刺透了祂的肋旁；我们尊贵的父亲若瑟求得了祂的身体；祂复活了，并且，如他们所说，三位教师看见祂被接到天上去了。}拉比肋未也为西默盎长老所说的话作了见证——祂被立定，使以色列中许多人跌倒和兴起，并成为一个反对的记号。\par

那时，教师狄达斯对全会众说：\Quote{如果这些人为耶稣所见证的一切事都成就了，那么这些事就是出于天主，在我们眼中不应感到惊奇。}会堂的首领、司祭和肋未人彼此说：\Quote{我们的法律如何记载呢？它说：\InnerQuote{祂的名当永远受赞颂：祂的居所在太阳之前常存，祂的宝座在月亮之前；地上万族都要因祂蒙福，万民都要事奉祂；远方的君王要来朝拜祂，尊敬祂。}}\par

\SourceRecord{Translated by Alexander Walker.；Ante-Nicene Fathers；Vol. 8.；Edited by Alexander Roberts, James Donaldson, and A. Cleveland Coxe.；Buffalo, NY: Christian Literature Publishing Co.,；1886.；<http://www.newadvent.org/fathers/08071c.htm>.}

% structure: p0001 source=h1 target=section reason=page-title

\section{\WorkTitle{尼苛德摩福音（第二部分，希腊文版）}}

\Emph{基督下降阴府。}\par

% structure: p0003 source=h2 target=subsection reason=relative-heading-rank; base=h2

\subsection{第一章（17）。}

约瑟说：\Quote{你们为什么惊讶耶稣复活了呢？但更奇妙的是，他不仅自己复活了，还使许多其他死人复活了，他们出现在耶路撒冷很多人面前。\BibleRef{玛 27:53}如果你们不认识其他人，至少\Person{西默盎}——他曾怀抱耶稣——以及他所复活的两个儿子，你们总该认识吧。因为我们不久前才埋葬了他们；但现在他们的坟墓是开的\Emph{且}空着，他们活着，住在\Place{阿黎玛特雅}。}于是他们派人去，发现坟墓是开的且空着。约瑟说：\Quote{我们到\Place{阿黎玛特雅}去找他们吧。}\par

于是司祭长\Person{亚纳斯}和\Person{盖法}，以及\Person{约瑟}、\Person{尼苛德摩}、\Person{加玛里耳}和与他们同去的人起身，前往\Place{阿黎玛特雅}，找到了约瑟所说的人。他们做了祈祷，彼此问候。然后他们与那些人一起回到\Place{耶路撒冷}，把他们带进会堂，关上大门，把犹太人的旧\Emph{盟约}放在中间；司祭长对他们说：\Quote{我们希望你们指着以色列的天主和\OriginalTerm{阿多乃}{Adonai}起誓，要你们说实话，你们是如何复活的，是谁使你们从死人中复活的。}\par

那些复活的人听了这话，就在自己脸上画了十字圣号，对司祭长说：\Quote{请给我们纸、墨水和笔。}他们便拿来了。于是他们坐下，这样写道：——\par

% structure: p0007 source=h2 target=subsection reason=relative-heading-rank; base=h2

\subsection{第二章（18）。}

主耶稣基督，世界的复活与生命，求你赐予我们恩宠，使我们能述说你的复活，以及你在\OriginalTerm{阴府}{Hades}中所行的奇迹。我们当时正与自创世以来所有长眠的人一同在阴府里。到了半夜时分，有一道光如同太阳升起来，照进这些黑暗的\Emph{区域}。我们都得了光照，看见了彼此。立刻，我们的先祖\Person{亚巴郎}与众圣祖和先知聚集在一起，同时充满喜乐，彼此说：这光是来自一个伟大的光源。当时在场的\Person{依撒意亚}先知说：\Quote{这光是来自父、子及圣神；我活着时曾预言过他，说：\BibleQuote{则步隆地和纳斐塔里地，那坐在黑暗中的百姓，见到了大光明。}\BibleRef{依 9:1}}\par

这时，又有一位从旷野来的苦修者来到他们中间；圣祖们问他：你是谁？他回答说：\Quote{我是\Person{若翰}，最后的先知，我曾预备天主的道路，\BibleRef{玛 3:3}并向民众宣讲悔改以得罪赦。\BibleRef{谷 1:4}天主子曾来到我这里；我远远看见他，就对民众说：\BibleQuote{请看，天主的羔羊，除免世罪者！}\BibleRef{若 1:29}我亲手在\Place{约旦河}里给他付了洗，看见圣神犹如鸽子降在他身上；我也听见天主圣父的声音这样说：\BibleQuote{这是我的爱子，我所喜悦的。}\BibleRef{路 3:22}因此，他也派我到你们这里来，宣布天主的独生子将要到来，凡信他的人必得救，不信的人必被定罪。因此，我对你们众人说：当你们看见他时，你们都要朝拜他；因为现在正是你们悔改的时候，为你们在虚浮的上界敬拜偶像所犯的罪，以及你们所犯的其他罪过；而此后不再有机会。}\par

% structure: p0010 source=h2 target=subsection reason=relative-heading-rank; base=h2

\subsection{第三章（19）。}

因此，当\Person{若翰}正在阴府里教导众人时，首先受造的始祖\Person{亚当}听见了，就对他儿子\Person{舍特}说：\Quote{我儿，我愿你在我要死的时候，告诉人类的祖先和先知我曾派你到什么地方去。}\Person{舍特}说：\Quote{先知和圣祖们，请听。有一次，当我父亲\Person{亚当}，首先受造的，即将面临死亡时，他派我紧挨着乐园之门向天主恳求，求他派一位天使领我到怜悯树那里，以便我取油来敷抹我父亲，使他从病中起来：这件事，我也确实做了。祈祷之后，一位上主的天使来到，对我说：\InnerQuote{舍特，你求什么呢？你求那能使病人痊愈的油，还是求这油所流出的树，为了你父亲的病？这事现在不能找到。所以，你回去告诉你父亲：从世界受造起，直到满了五千五百年，那时天主的独生子将降生成人来到世上；他将用这油敷抹他，使他复活；并用水和圣神洗净他以及从他而出的人，然后他将治好一切疾病；但现在这事是不可能的。}}\par

当圣祖们和先知们听了这些话，非常欢喜。\par

% structure: p0013 source=h2 target=subsection reason=relative-heading-rank; base=h2

\subsection{第四章（20）。}

当众人都这样欢喜时，\OriginalTerm{撒殚}{Satan}，黑暗的继承人，来对阴府说：\Quote{噢，吞噬一切且永不满足者，听我的话。在犹太人中有一个名叫耶稣的，自称是天主子；他本是个人，通过我们与犹太人合作，犹太人把他钉死了；现在他死了，准备好，我们好把他扣押在这里。因为我知道他是个人，我也听见他说过：\BibleQuote{我的心灵忧闷得要死。}\BibleRef{谷 15:34}他在与世人同住在上界时，也对我做了许多恶事。因为无论在哪里遇到我的仆人，他就迫害他们；凡是我使人成为驼背、瞎眼、瘸腿、癞病的，或诸如此类的，他一句话就治好了他们；许多我已经准备好要埋葬的人，他甚至用一句话就使他们复活了。}\par

阴府说：\Quote{这个\Emph{人}竟如此有能力，能用一句话做这些事吗？或者，如果他真这样，你能抵挡他吗？在我看来，如果他这样，就没有人能抵挡他。如果你说你听见他畏惧死亡，他这样说是在嘲弄你，在讥笑你，想用强手抓住你；祸哉，祸哉，你将永远受祸！}\par

撒殫說：\Quote{噢，吞噬一切、貪得無厭的陰府，你聽到我們共同的敵人，就如此害怕嗎？我並不害怕他，反而在猶太人中間行事，他們把他釘在十字架上，還拿苦膽調醋給他喝。\BibleRef{瑪 27:34} 那麼準備好，好當他來的時候，你能牢牢抓住他。}\par

陰府回答說：\Quote{黑暗的繼承者、毀滅之子、魔鬼，你剛才告訴我，那些你預備埋葬的許多人，被一句話就復活了。如果他已把別人從墳墓中救出，我們如何並用什麼能力抓住他？因為我不久前吞下一個死人，名叫拉匝祿；不久之後，一個活人用一句話就強行把他從我的腹中拉了出來：我認為那就是你所說的那位。因此，如果我們在這裡接待他，我怕恐怕連其餘的人也會有危險。因為，看，我從永恆所吞下的所有那些人，我察覺他們在動盪，我的肚子疼痛。而且拉匝祿被奪走，在我看來不是好兆頭：因為他不像一個死屍，而是像一隻鷹，從我裡面飛了出去；因為地如此突然地把他拋出。因此我也懇求你，為了你的利益也為了我，不要帶他到這裡來；因為我認為他來這裡是為了復活所有的死人。我告訴你：憑我們所處的黑暗，如果你帶他到這裡來，沒有一個死人會留在我這裡。}\par

% structure: p0018 source=h2 target=subsection reason=relative-heading-rank; base=h2

\subsection{第五章（21）。}

撒殫和陰府這樣彼此說話時，有一個像雷鳴的巨大聲音說：\Quote{諸位首領，你們要抬起門戶；永久的門戶，你們也要被高舉；榮耀的王將要進來。}陰府聽見，就對撒殫說：\Quote{你出去，如果你能，就抵擋他。}撒殫於是走到外面。然後陰府對他的魔鬼說：\Quote{牢牢地、堅固地鎖住銅門和鐵閂，注意我的門閂，站好次序，察看一切；因為如果他進來，禍患就會抓住我們。}\par

先祖們聽見這話，都開始責罵他，說：\Quote{噢，吞噬一切、貪得無厭的！打開，好讓榮耀的王進來。}先知達味說：\Quote{瞎眼的，你不知道嗎？我活在世上時曾預言這句話：\InnerQuote{諸位首領，你們要抬起門戶。}}依撒意亞說：\Quote{我藉著聖神預見這事，寫道：\InnerQuote{死人要復活，墳墓中的人要起來，地上的人要歡樂。}死亡，你的刺在哪裡？陰府，你的勝利在哪裡？}\BibleRef{歐 13:14}\par

然後又有聲音說：\Quote{抬起門戶。}陰府第二次聽見這聲音，好像假裝不知道，回答說：\Quote{這榮耀的王是誰？}主的使者說：\Quote{主大能而有力，主在戰場上有能。}立刻隨著這些話，銅門破碎，鐵閂折斷，所有被捆綁的死人從監獄中出來，我們也和他們一同出來。榮耀的王以人的形狀進來，陰府所有的黑暗地方都被照亮了。\par

% structure: p0022 source=h2 target=subsection reason=relative-heading-rank; base=h2

\subsection{第六章（22）。}

陰府立刻喊叫說：\Quote{我們被征服了：禍哉我們！但你是誰，有這樣的能力和權柄？你是什麼，來到這裡而無罪，看起來微小卻大有能力，卑微卻崇高，既是奴僕又是主人，既是士兵又是君王，掌管死人和活人？你被釘在十字架上，被放在墳墓裡；現在你自由了，毀滅了我們所有的權力。那麼你就是耶穌，就是首領撒殫告訴我們的那位，說你要藉著十字架和死亡來繼承整個世界嗎？}\par

然後榮耀的王抓住首領撒殫的頭，把他交給他的天使，說：\Quote{用鐵鏈捆綁他的手、腳、脖子和口。}然後把他交給陰府，說：\Quote{帶走他，好好看管，直到我第二次顯現。}\par

% structure: p0025 source=h2 target=subsection reason=relative-heading-rank; base=h2

\subsection{第七章（23）。}

陰府接收撒殫，對他說：\Quote{貝耳則步，火與刑罰的繼承者，聖人的仇敵，你出於什麼必要，竟使榮耀的王被釘十字架，好讓他來到這裡，奪去\Emph{我們的權力}？轉身看看，沒有一個死人留在我這裡，相反，你通過知善惡樹所獲得的一切，都藉著十字架上的木頭失去了：你所有的喜樂都變成了憂傷；你想殺死榮耀的王，卻殺死了你自己。因為，既然我接收你來保護你，你將通過經驗知道我將對你行多少惡事。哦，大魔鬼，死亡的開端，罪惡的根源，一切邪惡的終結，你在耶穌身上發現了什麼惡事，竟要圖謀毀滅他？你怎麼敢做這樣邪惡的事？你怎麼忙於把這樣一個人帶進這黑暗，藉著他你已經被剝奪了所有從永恆死去的人？}\par

% structure: p0027 source=h2 target=subsection reason=relative-heading-rank; base=h2

\subsection{第八章（24）。}

陰府這樣對撒殫說話時，榮耀的王伸出他的右手，抓住我們的先祖亞當，把他舉起來。然後也轉向其餘的人，說：\Quote{所有因他所摸的樹而死的人，都跟我來；因為，看，我藉著十字架上的木頭再次使你們全部復活。}於是他把他們全都帶出來，我們的先祖亞當似乎充滿喜樂，說：\Quote{主啊，我感謝你的威嚴，因為你把我從最低的陰府中提上來。}同樣，所有的先知和聖人都說：\Quote{我們感謝你，基督，世界的救主，因為你把我們的生命從毀滅中提上來。}\par

他們這樣說完之後，救主用十字架的記號祝福亞當的額頭，也同樣祝福了族長、先知、殉道者和先祖們；他帶著他們，從陰府中躍升出來。他行走時，聖父們陪同他，唱著讚歌說：\Quote{奉主名而來的，當受讚美：亞肋路亞；願眾聖人的榮耀都歸於他。}\par

% structure: p0030 source=h2 target=subsection reason=relative-heading-rank; base=h2

\subsection{第九章（25）。}

于是，他起身前往\Place{乐园}，拉住我们始祖\Person{亚当}的手，将他与所有义人交给总领天使\Person{弥额尔}。当他们进入乐园门口时，有两位老人迎接他们，圣祖们问他们说：\Quote{你们是谁？你们没有见过死亡，没有下到阴间，却带着肉身和灵魂居住在乐园？}其中一位回答说：\Quote{我是\Person{哈诺客}，曾中悦天主，被祂提到这里；这位是特士布的\Person{厄里亚}。我们也要活到世界终结，那时我们要被天主派遣去对抗假基督，被他杀死，三天后复活，被提到云中，迎接主。}\par

% structure: p0032 source=h2 target=subsection reason=relative-heading-rank; base=h2

\subsection{第十章（26）}

正当他们这样说时，又来了一个谦卑的人，肩上背着一个十字架。圣祖们对他说：\Quote{你是谁？看起来像个强盗。你肩上背的十字架是什么？}他回答说：\Quote{我，如你们所说，在世上曾是个强盗和贼，为此犹太人抓住我，将我连同我们的主耶稣基督一起交给十字架之死。当祂悬在十字架上时，我看见了所行的奇迹，就相信了祂，并向祂恳求说：\InnerQuote{主，当你成为君王时，请不要忘记我。}祂立刻对我说：\InnerQuote{阿们，阿们：我今天告诉你，你将与我同在乐园里。}因此我背着十字架来到乐园；我遇见总领天使\Person{弥额尔}，对他说：\InnerQuote{我们的主耶稣，那位被钉十字架者，派我来到这里；所以请带我到伊甸园门口。}火焰剑看见十字架的记号，就为我打开了，我便进去了。于是总领天使对我说：\InnerQuote{稍等片刻，因为人类的始祖亚当也带着义人前来，他们也要进去。}现在，我看见你们，就出来迎接你们。}\par

圣徒们听到这话，都大声呼喊说：\Quote{我们的主何其伟大，祂的德能何其伟大！}\par

% structure: p0035 source=h2 target=subsection reason=relative-heading-rank; base=h2

\subsection{第十一章（27）}

所有这一切我们都亲眼看见、亲耳听见；我们两兄弟，被总领天使\Person{弥额尔}派遣，受命宣告主的复活，但先要去约旦河受洗。我们也去了那里，与其他复活了的死者一起受了洗。此后我们来到耶路撒冷，庆祝了复活的逾越节。但现在我们要离开了，不能留在这里。愿天主、圣父的慈爱，我们的主耶稣基督的恩宠，圣神的共融，与你们众人同在。\BibleRef{格后 13:15}\par

写完这些并封好卷轴后，他们将一半交给了大司祭，一半交给了\Person{若瑟}和\Person{尼苛德摩}。他们便立刻消失了：归于我们的主耶稣基督的光荣。阿们。\par

\SourceRecord{Translated by Alexander Walker.；Ante-Nicene Fathers；Vol. 8.；Edited by Alexander Roberts, James Donaldson, and A. Cleveland Coxe.；Buffalo, NY: Christian Literature Publishing Co.,；1886.；<http://www.newadvent.org/fathers/08072a.htm>.}

% structure: p0001 source=h1 target=section reason=page-title

\section{尼苛德摩福音（第二部分，第一拉丁文本）}

\Emph{基督下降阴府}\par

% structure: p0003 source=h2 target=subsection reason=relative-heading-rank; base=h2

\subsection{第一章（17）}

\Person{若瑟}站起来对\Person{亚纳斯}和\Person{盖法}说：你们确实有理由惊讶，因为你们听说耶稣从死者中复活，升到天上。但更令人惊奇的是，他并非唯一从死者中复活的人；他还使许多其他死者从坟墓中活过来，被\Place{耶路撒冷}许多人看见。现在请听我说，我们都知道那位伟大的司铎圣\Person{西默盎}，他在圣殿中怀抱婴孩耶稣。西默盎本人有两个儿子，是亲兄弟；我们都在他们安眠和埋葬时在场。所以，去吧，去看看他们的坟墓：它们是敞开的，因为他们已经复活了；看，他们就在\Place{阿黎玛特雅}城里，一起祈祷度日。确实，他们发出呼喊声，却不与任何人说话，像死人一样沉默。但来，我们去他们那里；我们以一切尊敬和礼仪带他们到我们这里。如果我们命令他们，也许他们会向我们讲述他们复活的奥秘。\par

听到这个，他们都喜乐了。\Person{亚纳斯}和\Person{盖法}、\Person{尼苛德摩}、\Person{若瑟}和\Person{加玛里耳}去了，却没有在坟墓里找到他们；但走进\Place{阿黎玛特雅}城，他们在那里找到了他们，屈膝祈祷。他们亲吻了他们，以一切敬畏和敬畏天主的心，带他们到\Place{耶路撒冷}的会堂。关上门，举起上主的法律书，放在他们手中，以\OriginalTerm{阿多乃}{Adonai}天主和以色列的天主——祂通过法律和先知向我们祖先说话——命令他们，说：\Quote{你们相信是耶稣使你们从死者中复活吗？告诉我们你们是如何从死者中复活的。}\par

\Person{卡里努斯}和\Person{娄基乌斯}听到这个命令，身体颤抖，心中不安而叹息。他们一同仰望上天，用手指在舌头上画了十字圣号，立刻一起说：\Quote{给我们每人一张纸，让我们写下我们所看见和听见的。}他们就给了他们。他们坐下，各自写道：\par

% structure: p0007 source=h2 target=subsection reason=relative-heading-rank; base=h2

\subsection{第二章（18）}

\Person{主耶稣基督}，死者的复活和生命，求祢允许我们通过祢十字架的死亡讲述奥秘，因为我们受祢命令所约束。因为祢命令祢的仆人不要向任何人讲述祢在阴府中所行的神圣威严的奥秘。当我们与所有祖先一同躺在深渊的黑暗之中时，忽然出现了金色的太阳热气和紫色的君王之光，照耀着我们。立刻，整个人类的祖先连同所有的圣祖和先知都欢呼说：\Quote{那光是永恒之光的泉源，祂应许将同永恒之光传给我们。}\Person{依撒意亚}呼喊说：\Quote{这是圣父的光，天主子，正如我活在世上时所预言的：}\BibleQuote{在约但河对岸的则步隆地和纳斐塔里地，外邦人的加里肋亚，坐在黑暗中的百姓看见了皓光；那些住在死亡阴影地区的人，有光照耀他们。}现在这光来了，照耀了我们这些坐在死亡中的人。\par

当我们都在照耀我们的光中欢呼时，我们的父亲\Person{西默盎}来到我们这里，欢呼说：\Quote{请光荣主耶稣基督，天主子；因为我在圣殿中亲手抱过祂还是婴孩的时候，受圣神感动，我向祂承认说：}\BibleQuote{现在我的眼睛看见了祢的救恩，就是祢在万民面前所预备的，为启示外邦人的光，和祢子民以色列的荣耀。}当众人听到这话时，所有圣徒的群众更加欢腾。\par

此后，有一位像是旷野的居民上来；大家问他：你是谁？他回答说：\Quote{我是\Person{若翰}，至高者的声音和先知，走在祂来临前面预备祂的道路，将救恩的知识赐给祂的子民，以赦免他们的罪。我看见祂向我走来，受圣神感动，我说：\InnerQuote{看，天主的羔羊！看，除免世罪者！}我在约但河里给祂付了洗，我看见圣神像鸽子降临在祂身上；我听见从天上来的声音说：\InnerQuote{这是我的爱子，我所喜悦的。}现在我已经走在祂前面，降临来向你们宣告，那升起的圣子即将来临，从高天来访问我们这些坐在黑暗和死亡阴影中的人。}\par

% structure: p0011 source=h2 target=subsection reason=relative-heading-rank; base=h2

\subsection{第三章（19）}

当第一个受造者，父亲\Person{亚当}，听到耶稣在\Place{约但河}受洗时，他向他儿子\Person{舍特}喊叫：\Quote{告诉你的儿子们，圣祖和先知们，所有你从我派你到乐园门口祈求天主派遣祂的天使赐给你\OriginalTerm{仁慈之树}{tree of mercy}的油时，从总领天使\Person{弥额尔}那里听到的一切，为的是当我生病时用那油膏我的身体。}于是\Person{舍特}走近圣祖和先知们，说：\Quote{当我\Person{舍特}在乐园门口向主祈祷时，看，上主的天使\Person{弥额尔}向我显现，说：\InnerQuote{我是上主派来给你的。我掌管人类。我告诉你\Person{舍特}，不要因油膏你父亲\Person{亚当}以解他身体疼痛的眼泪祈祷和恳求，因为你决不会得到它，除非在末后的日子和时期，直到五千年五百年满了：那时天主最爱的圣子将来到地上，再次复活\Person{亚当}的身体和死者的身体；祂来到时，将在\Place{约但河}受洗。当祂从\Place{约但河}的水中出来时，祂将用祂\OriginalTerm{仁慈之油}{oil of mercy}膏所有相信祂的人；那\OriginalTerm{仁慈之油}{oil of mercy}将为那些从水和圣神重生得永生的人而存留。然后，基督耶稣，天主最爱的圣子，将降临地上，引领我们的父亲\Person{亚当}进入乐园，到达\OriginalTerm{仁慈之树}{tree of mercy}。}}\par

当\Person{舍特}讲了这一切，所有圣祖和先知都大大欢欣。\par

% structure: p0014 source=h2 target=subsection reason=relative-heading-rank; base=h2

\subsection{第四章（20）}

当众圣人都在欢欣雀跃时，看啊，死亡的君王和首领撒殚对哈德斯说：你准备迎接耶稣吧，祂自称为天主子，却是一个惧怕死亡的人，说：\Quote{我的心灵忧闷得要死。}祂曾多次抵挡我，加害于我；许多被我弄瞎、瘸腿、耳聋、患麻风病、附魔的人，祂用一句话就治好了；而那些被我带到你这里的死人，祂却从你手中夺走了。\par

哈德斯回答撒殚首领说：那一位是谁，竟如此强大，而祂不过是一个惧怕死亡的人？因为地上所有有权柄的人都被我的权柄所制服，而你用你的权柄使他们服从。那么，如果你是有权柄的，那个耶稣又是怎样的人呢？祂虽然怕死，却能抵挡你的权柄？如果祂在人性中如此强大，我实在告诉你，祂在神性中全能，无人能抵抗祂的权柄。当祂说自己怕死时，祂是要抓住你，而你将永永远远有祸了。塔尔塔罗斯首领撒殚回答说：你为什么怀疑并害怕接待这个耶稣，祂是你的敌人，也是我的敌人？因为我曾试探祂，我煽动我古时的百姓犹太人，以仇恨和愤怒反对祂；我磨尖了长矛要刺祂；我调和了苦胆和醋给祂喝；我预备了木头来钉死祂，预备了钉子来刺穿祂；祂的死期近了，好叫我把祂带到你面前，使祂服从你和我。\par

塔尔塔罗斯回答说：你曾告诉我，正是祂从我这里夺走了死人。现在有许多被我拘禁在这里的人，当他们在地上活着时，曾从我这里取走死人，不是靠他们自己的能力，而是靠虔诚的祈祷，他们的全能天主把他们从我这里夺走了。那个耶稣是谁，竟能不用祈祷，一句话就把死人从我这里带走？也许祂就是用命令的话使拉匝禄复活的同一人，那拉匝禄在臭气和腐烂中过了四天，我把他当作死人拘留着。死亡的君王撒殚回答说：那个耶稣就是同一人。哈德斯听了这话就对他说：我指着你的权柄和我的权柄起誓，你不要把祂带到我这里来。因为那时，当我听到祂命令的话时，我因恐惧和惊慌而战栗，我所有的官员也同时与我一同惊慌失措。我们也不能留住那个拉匝禄；但他像鹰一样抖擞，迅速敏捷地从我们这里跃出，走了出去，那曾经容纳拉匝禄尸体的同一地面立刻把他活着交了出来。所以现在，我知道那能做这些事的人就是天主，祂在权柄上坚强，在人性中有大能，祂是人类救主。但如果你把祂带到我这里，凡被关在这残酷监狱中、被他们的罪孽用不能松开的锁链捆绑的人，祂都要释放，并带他们进入祂神性的永生，直到永远。\par

% structure: p0018 source=h2 target=subsection reason=relative-heading-rank; base=h2

\subsection{第五章（21）}

撒殚首领和哈德斯这样轮番对话时，忽然有声音如同雷鸣，并有灵体的呼喊声：\Quote{众城门哪，抬起你们的头来；永久的门户啊，你们要被举起；光荣的君王要进来。}哈德斯听见这话，对撒殚首领说：你离开我，滚出我的领域；你若是有能的战士，就去与光荣的君王作战。但你和祂有什么关系？哈德斯就把撒殚赶出了他的领域。哈德斯对他邪恶的官员说：关上残忍的铜门，插上铁闩，勇敢地抵抗，免得我们这些被掳的人反被\Emph{祂}掳去。\par

众圣人听见这话，就用责备的声音对哈德斯说：\Quote{敞开你的城门，让光荣的君王进来。}达味喊道说：我在地上活着的时候，不是曾向你们预言说：\Quote{但愿人因上主的仁慈和祂向人子所做的奇事，都称谢上主，因为祂击碎了铜门，砍断了铁闩；祂从他们的邪路上救拔了他们？}之后，同样地，依撒意亚说：我在地上活着的时候，不是曾向你们预言说：\Quote{死人必复活，在坟墓里的必兴起，在地上的人要欢欣，因为从上主来的露水是他们的健康？}我又说：\Quote{死亡啊，你的刺在哪里？哈德斯啊，你的胜利在哪里？}\par

众圣人听了依撒意亚这话，就对哈德斯说：\Quote{敞开你的城门。如今你已被征服，你将软弱无力。}有大声音如同雷鸣，说：\Quote{众城门哪，抬起你们的头来；阴间的门户啊，你们要被举起；光荣的君王要进来。}哈德斯见他们两次这样呼喊，就装作不认识说：光荣的君王是谁？达味回答哈德斯说：我认出这呼喊的话，因为我曾借着祂的神预言了同样的话。现在，我把我上面所说的告诉你：\Quote{上主大能有力，上主作战有力；祂就是光荣的君王。}上主亲自从天垂看大地，\Quote{要听被囚之人的叹息，要释放被定死罪的人。}现在，最污秽最可憎的哈德斯啊，敞开你的城门，让光荣的君王进来。达味这样说时，有一位荣耀的主以人的形像来到哈德斯那里，照亮了永久的黑暗，挣断了那不可解除的锁链；无敌大能的援助临到了我们这些坐在罪孽深重之黑暗中和死于罪恶之阴影里的人。\par

% structure: p0022 source=h2 target=subsection reason=relative-heading-rank; base=h2

\subsection{第六章（22）}

当阴府与死亡，以及他们不虔敬的官员、残忍的仆役看见这事，他们在自己领域内察觉如此大光明的灿烂，又看见基督突然出现在他们的居所，就战栗起来，喊叫说：\Quote{我们被你战胜了。你是谁，竟向主指摘我们的混乱？你是谁，未被腐朽毁灭，你威能的纯净证明，竟愤怒地谴责我们的权势？你是谁，如此伟大又微小，卑下又崇高，士兵又统帅，以奴隶形象出现的奇妙战士，且是死而复生的光荣的君王，被十字架所承载的受宰者？你，曾躺在坟墓中死去，却活着下到我们这里；在你的死亡中，一切受造物战栗，众星一同震动；如今你在死者中得了自由，扰乱我们的军团。你是谁，释放那些被原罪捆绑的俘虏，召叫他们恢复原有的自由？你是谁，将神圣、灿烂、光明的光辉倾注在那些被自己罪恶的黑暗所蒙蔽的人身上？}\par

同样，所有恶魔的军团，从他们可怕的倾覆中，被同样的恐惧所惊骇，喊叫说：\Quote{耶稣，你是从哪里来的，一个人如此强大、荣耀辉煌、卓越无瑕、清白无罪？因为那个一直臣服于我们直到现在的尘世，曾为我们所用纳税，从未给我们送来这样一个死人，从未为下面的权势预定如此礼物。那么你是谁，竟然如此无畏地进入我们的边界，不仅不怕我们的惩罚，反而设法将所有人从我们的锁链中夺走？也许你就是我们的魔王撒旦所说的那个耶稣，他说你将因十字架之死注定要获得全世界的统治权。}\par

那时，光荣的君王以其威能践踏死亡，擒拿魔王撒旦，将他交于阴府之权下，并将亚当引向自己的光明。\par

% structure: p0026 source=h2 target=subsection reason=relative-heading-rank; base=h2

\subsection{第七章（23）}

于是阴府接受魔王撒旦，以强烈的斥责对他说：\Quote{啊，毁灭的王子，灭绝的首领，贝耳则步，天使的嘲笑对象，被义人唾弃的，你为什么想做这事？你想把光荣的君王钉在十字架上，在他的死中你曾答应我们如此巨大的战利品吗？你像个愚人，不知道自己在做什么。看，那个耶稣以其天主性的光辉驱散了死亡的一切黑暗，他闯入了我们牢狱的最深坚固之处，带出了俘虏，释放了被捆绑的人。所有曾在我们折磨下呻吟的人都辱骂我们，因他们的祈祷，我们的统治被攻占，我们的领域被征服，现在没有哪个人类种族对我们有任何敬畏。此外，我们还受到死人的严重威胁，他们从未向我们傲慢，也从未像俘虏一样快乐过。啊，魔王撒旦，所有不虔敬的恶人与叛徒之父，你为什么想做这事？那些从起初直到现在对救恩和生命绝望的人，这里再也听不到他们惯常的嚎叫；他们的呻吟也不复回响，他们的脸上寻不到泪痕。啊，魔王撒旦，地下区域的钥匙持有者，你因悖逆之树和丧失乐园所获得的一切财富，如今因十字架之树而尽失，你的一切喜乐都已消亡。当你挂起那位基督耶稣——光荣的君王时，你是在反对自己和反对我。从现在起，你将知道你要在我永久的看守中忍受怎样的永恒苦刑和无限惩罚。啊，魔王撒旦，死亡之作者，一切骄傲之源，你本该先查问那个耶稣的恶行。在你认为他毫无过犯的情况下，你为什么无缘无故地胆敢不义地将他钉死？你为什么将无辜和正义的人带到我们的区域，却失去了全世界的罪人、不虔敬和不义的人？}\par

阴府对魔王撒旦说了这些话后，光荣的君王就对阴府说：\Quote{魔王撒旦将永远在你的权下，代替亚当和他的子孙——我的义人们。}\par

% structure: p0029 source=h2 target=subsection reason=relative-heading-rank; base=h2

\subsection{第八章（24）}

那时，主伸出手来说：\Quote{到我这里来，我所有的圣人，你们拥有我的肖像和模样。你们曾因树、魔鬼和死亡而被定罪，现在看，魔鬼和死亡因树而被定罪。}众圣人立即聚集在主的手下。主握住亚当的右手，对他说：\Quote{愿平安归于你，以及你所有的子孙——我的义人们！}亚当俯伏在主膝前，含泪恳求，高声说：\BibleQuote{上主，我要宣扬你，因为你提拔了我，没有让我的仇敌因我而喜乐。上主天主，我曾呼求你，你医治了我。上主，你从下面的权势中领出了我的灵魂，你拯救了我，使我免于下到深坑。上主的众圣人，请歌颂上主，称谢他的圣名；因为他的震怒中有愤怒，他的善意中有生命。}同样，所有天主的圣人也都俯伏在主脚前，异口同声地说：\Quote{你来了，世界的救赎者！正如你藉法律和你的先知所预言的，你已藉你的作为成就了。你藉你的十字架救赎了活人，并藉十字架之死下到我们这里，为了藉你的威能把我们从下面的权势和死亡中拯救出来。上主，正如你在天上设立了你的荣耀的称号，并在世上竖立了你的十字架作为救赎的称号，同样，上主，请在阴府设立你十字架胜利的记号，使死亡不再作王。}\par

那时，主伸出手，在亚当和所有他的圣人身上画了十字圣号；然后握住亚当的右手，从下面的权势中上升；所有圣人都跟随他。接着，圣达味高声喊道：\BibleQuote{请向上主高唱新歌，因为他行了奇事；他的右手和他的圣臂为他带来了救恩。上主彰显了他的救恩，他在外邦人眼前揭示了正义。}全体圣众回答说：\Quote{这是所有圣人的荣耀。阿们，阿肋路亚。}\par

此后，\Person{哈巴谷}先知呼喊说：\BibleQuote{祢出来救恩祢的百姓，救恩祢的选民。}\BibleRef{哈 3：13}众圣人都回答说：\BibleQuote{奉主名而来的当受赞颂；天主是主，祂光照了我们。阿们，阿肋路亚。}同样，\Person{米该亚}先知也呼喊说：\BibleQuote{上主，有谁能像祢，赦免罪愆，宽宥过犯？现在祢为了一个\Emph{反对我们}的证言而抑制祢的愤怒，因为祢喜爱仁慈。祢再回转，怜悯我们，赦免我们的一切罪愆；祢将我们的一切罪过投于深海，正如祢古时向我们的祖先所起誓的。}\BibleRef{米 7：18}众圣人都回答说：\BibleQuote{这是我们的天主，直到永远，世世无穷；祂将永远引导我们。阿们，阿肋路亚。}同样，所有先知，引用关于祂赞美的神圣\Emph{著作}，以及所有圣人呼喊阿们、阿肋路亚，都跟随了主。\par

% structure: p0033 source=h2 target=subsection reason=relative-heading-rank; base=h2

\subsection{第九章（25）}

主握着\Person{亚当}的手，把他交给\Person{总领天使弥额尔}；众圣人都跟随\Person{总领天使弥额尔}，他带领他们全都进入乐园的荣耀恩宠中。那里有两位年高迈古的人迎接他们。圣人们问他们：你们是谁，还没有和我们一同在阴府中死去，却带着身体被安置在乐园里？其中一个回答说：我是\Person{哈诺客}，因主的话被迁到这里；与我同在的是\Person{提市贝人厄里亚}，他被火马车接去。我们至今在这里未尝死味，但被保留到假基督来临，凭神圣的奇迹奇事与他争战，并在\Place{耶路撒冷}被他杀死，三天半后又被活活提到云中。\par

% structure: p0035 source=h2 target=subsection reason=relative-heading-rank; base=h2

\subsection{第十章（26）}

当圣\Person{哈诺客}和\Person{厄里亚}这样说时，看，另有一个极其可怜的人走上来，肩上背着十字架的标记。众圣人看见他，都问他说：你是谁？因为你的相貌像个强盗。你肩上背的是什么标记？他回答他们说：\Quote{你们说得对，我本来是个强盗，在地上作尽各种恶事。犹太人把我与耶稣一同钉在十字架上；我看见受造之物通过被钉耶稣的十字架所行的奇迹，就相信祂是万有的创造者，全能之王；我恳求祂说：\InnerQuote{主，当祢进入祢王国时，请记念我。}祂立刻接纳我的恳求，对我说：\InnerQuote{阿们，我告诉你：今天你就要同我在乐园里。}\BibleRef{路 23：42}祂把这个十字架的标记给了我，说：\InnerQuote{带着这个走进乐园；如果乐园的守护天使不让你进去，你就给他看十字架的标记，并对他说：\InnerQuote{现在被钉十字架的耶稣基督，天主子，派我来了。}}我这样做了，将这一切告诉了乐园的守护天使。他听见后，立刻打开门，领我进去，把我安置在乐园的右边，说：\InnerQuote{看，稍等片刻，全人类之父亚当就要带着他所有圣洁公义的子女进来，在基督——被钉十字架的主——升天的胜利和荣耀之后。}}强盗的这番话，所有圣祖和先知都异口同声说：\Quote{主全能的天主，永远恩惠之父，仁慈之父，祢是应当赞美的！祢赐予祢的罪人如此恩宠，使他们回到乐园的恩宠和祢丰美的草场中；因为这是最确实的属灵生命。阿们，阿们。}\par

% structure: p0037 source=h2 target=subsection reason=relative-heading-rank; base=h2

\subsection{第十一章（27）}

这些是神圣的奥秘，我们\Person{卡里努斯}和\Person{路济乌斯}所看见听见的。天主其他的奥秘，我们不得多讲，因为总领天使弥额尔命令我们，说：\Quote{你们要到\Place{耶路撒冷}与你们的弟兄同去，持续祈祷，呼喊并光荣主耶稣基督的复活，祂使你们与祂一同从死者中复活。不与任何人说话，要像哑巴一样坐着，直到主亲自允许你们讲述祂天主性的奥秘的时候。}总领天使弥额尔又命令我们渡过\Place{约但河}，进入一个富饶肥沃的地方，那里有许多同我们一起复活的人，作为主基督复活的证据；因为我们这些从死者中复活的人，只被允许三天在\Place{耶路撒冷}与活着的亲属一起庆祝主的逾越节，作为主基督复活的证据；我们在圣\Place{约但河}中领受了圣洗圣事，每人接受了白衣。三天后，我们庆祝了主的逾越节，所有同我们一起复活的人都被提到云中，被带到约但河对岸，再也不被任何人看见。但我们被告知要留在\Place{阿黎玛特雅}城祈祷。\par

这些是主命令我们告诉你们的事。你们要赞美和宣认祂，并悔改，使祂怜悯你们。愿同一主耶稣基督，我们众人的救主，赐给你们平安！阿们。\par

他们各自在分开的纸上写完一切后，站了起来。\Person{卡里努斯}把所写的交给\Person{亚纳斯}、\Person{盖法}和\Person{加玛里耳}；同样，\Person{路济乌斯}把所写的交给\Person{尼苛德摩}和\Person{若瑟}。他们忽然变了形貌，变得极其洁白，就不再被看见。他们的著作被发现完全一致，不多一个字母也不少一个字母。\par

犹太人的整个会堂，听见\Person{卡里努斯}和\Person{路济乌斯}这些奇妙的话语，彼此说：\Quote{这些事确实都是主所做的，愿主永受赞颂，直到永远。阿们。}他们都极其焦虑地出去，恐惧战栗地捶着胸，各自回家去了。\par

犹太人在会堂里所说的这一切，\Person{尼苛德摩}和\Person{若瑟}立即报告给总督。\Person{比拉多}亲自写下了犹太人对耶稣所做所说的一切，并将所有的话放在他总督府的公共记录中。\par

% structure: p0043 source=h2 target=subsection reason=relative-heading-rank; base=h2

\subsection{第十二章（28）}

此后，彼拉多进入犹太人的圣殿，召集了所有的大司祭、学者、经师和法学士，同他们一起进入圣殿的内殿，命令关闭所有门户，对他们说：\Quote{我们听说你们在这圣殿里有一部伟大的经卷汇编；因此我要求你们将它呈现在我们面前。}当四位官员将那部镶有黄金和宝石的经卷汇编带来时，彼拉多对众人说：\Quote{我因你们祖先的天主——祂命令你们在这圣地建造圣殿——而恳求你们，不要向我隐瞒真相。你们都知道那部经卷汇编中写着什么；但现在你们要说，你们是否在经卷中发现你们钉死的耶稣就是要来为人类救恩的、应来的圣子，以及祂应在多少个季节轮回后来临。请向我说明，你们钉死祂是出于无知，还是明知故犯。}\par

被这样恳求后，亚纳斯和盖法命令所有同他们一起的人离开内殿；他们自己关闭了圣殿和内殿的所有门户，对彼拉多说：\Quote{我们被你恳求，贤明的审判官，因这座圣殿的建筑，要向你说明真相和这件\Emph{事}的清楚报告。我们在不知耶稣是圣子的情况下钉死了祂，以为祂是靠某种法术行奇迹；此后，我们在这圣殿里召开了一次大集会。我们在彼此商议耶稣所行奇迹的征兆时，发现我们民族中有许多见证人，他们说看见耶稣在受死之后还活着，并且升入了高天。我们又看见两个见证人，是耶稣从死者中复活的，他们告诉我们耶稣在死者中所行的许多奇妙之事，这些我们都记录在手中。我们的惯例是，每年在集会之前，打开那部神圣的经卷汇编，寻求天主的见证。我们在\WorkTitle{七十士译本}的第一卷书中发现，总领天使弥额尔曾向原祖亚当的第三子谈起五千五百年，那时基督、天主最钟爱的圣子将从天降下；由此我们考虑，祂或许是以色列的天主，祂曾对梅瑟说：\BibleQuote{\BibleRef{出 25:10}，你要为你制造约柜，长二肘半，宽一肘半，高一肘半。}在这五肘半中，我们从旧约约柜的结构中理解和认出，在五千五百年中，耶稣基督将要在身体的约柜中来临；我们找到了祂是以色列的天主、圣子。因为在祂受难后，我们这些大司祭因祂所发生的征兆而惊异，打开这部经卷汇编，搜寻所有世代，直到若瑟的世代，计算出基督的母亲玛利亚出自达味的后裔；我们发现，从天主创造天地和原祖到洪水，是二千二百一十二年；从洪水到巴贝耳塔建造，是五百三十一年；从塔的建造到亚巴郎，是六百零六年；从亚巴郎到以色列子民出离埃及，是四百七十年；从以色列子民出埃及到圣殿建造，是五百一十一年；从圣殿建造到圣殿毁坏，是四百六十四年。我们在\WorkTitle{厄斯德拉书}中找到了这些。经过搜索，我们发现从圣殿被焚到基督来临和祂诞生，是六百三十六年；总共是五千五百年，正如我们发现总领天使弥额尔向原祖亚当的第三子舍特所预言的，在五千五百年的末期，基督、圣子将降临。直到如今，我们没有告诉任何人，以免在我们的集会中引起纷争。现在，贤明的审判官，你因这部神圣的天主见证之书恳求我们，我们向你明示这些事。现在我们恳求你，以你的生命和健康，请不要在耶路撒冷向任何人透露这些话。}\par

% structure: p0046 source=h2 target=subsection reason=relative-heading-rank; base=h2

\subsection{第十三章（29）}

彼拉多听了亚纳斯和盖法的这些话，把它们全部记入我们主和救主的卷宗，存放在他的总督府档案中，并给罗马城的皇帝克劳狄写了一封信，说：\par

\Quote{本丢·彼拉多致其皇帝克劳狄，问安。最近发生的事，我自己也已证实，犹太人出于嫉妒，用残忍的刑罚惩罚了自己和他们的后代。简言之，当他们的祖先得到应许，他们的天主将从天上派遣祂的圣者给他们，祂理当被称为他们的君王，并应许祂将藉一童贞女降临人间；因此，当我任总督时，祂来到犹太，他们见祂使瞎子看见、洁净癞病人、治好瘫痪者、驱逐魔鬼、甚至复活死人、命令风浪、在海浪上干脚行走，并行了其他许多奇迹；所有犹太民众都说祂是圣子，大司祭们却对祂产生嫉妒，捉住祂，把祂交给我；他们连篇谎言说祂是个巫师，行事违反他们的法律。}\par

\Quote{我当时信以为真，便将祂交给他们鞭打，任凭他们处置。他们钉死了祂，埋葬后派兵看守。但第三天，在我士兵看守之时，祂复活了。犹太人的恶行如此昭彰，他们竟贿赂我的士兵，说：\InnerQuote{你们就说祂的门徒偷走了祂的身体。}但他们收了钱后，却不能隐瞒所发生的事；因为他们既见证祂已复活，又见证他们看见了祂，并承认从犹太人那里收了钱。}\par

\Quote{因此我做了这件事，免得有人编造不同的虚假说法，也免得你认为犹太人的谎言可信。}\par

\SourceRecord{Translated by Alexander Walker.；Ante-Nicene Fathers；Vol. 8.；Edited by Alexander Roberts, James Donaldson, and A. Cleveland Coxe.；Buffalo, NY: Christian Literature Publishing Co.,；1886.；<http://www.newadvent.org/fathers/08072b.htm>.}

% structure: p0001 source=h1 target=section reason=page-title

\section{\WorkTitle{尼苛德摩福音（第二部，第二拉丁文本）}}

% structure: p0002 source=h2 target=subsection reason=relative-heading-rank; base=h2

\subsection{\Emph{第一章（17）}}

那时，\Person{辣彼阿达斯}、\Person{辣彼斐乃斯}和\Person{辣彼厄基雅斯}——这三位从\Place{加里肋亚}来的人，曾作证看见耶稣被提到天上——在犹太首领的群众中站起来，在司祭和肋未人面前说（这些人已被召聚到主的议会中）：\Quote{当我们从\Place{加里肋亚}来的时候，我们在\Place{约但河}遇见了一大群人，是些早已死去之父辈。我们看见\Person{卡里诺斯}和\Person{路基乌斯}也在他们中间。他们走到我们跟前，我们彼此亲吻，因为他们是我们亲爱的朋友。我们问他们说：\InnerQuote{朋友们、兄弟们，请告诉我们，这是什么生命的气息和肉体？与你同行的那些人是谁？你们这些早已死去的人，怎么还留在身体内？}}\par

他们回答说：\Quote{我们与基督一同从阴府复活了，祂使我们从死者中复活。由此你们可以知道：死亡和黑暗的门已被摧毁，圣人们的灵魂已被从那里领出，并已与主基督一起升到天上。主亲自命令我们，在指定的时间内行走在\Place{约但河}岸和山岭间；但不是向每个人显现，也不是向每个人说话，只向祂允许我们显现的人说话。刚才若不是圣神允许我们，我们既不能向你们说话，也不能向你们显现。}\par

议会中所有的群众听了这话，都充满恐惧和战栗，想知道这些加里肋亚人所见证的事是否真的发生了。然后\Person{盖法}和\Person{亚纳斯}对议会说：\Quote{这些人的见证，从头到尾必须尽快完全查明：如果发现\Person{卡里诺斯}和\Person{路基乌斯}真的还活在身体内，并且我们能亲眼看见他们，那么他们的见证就全然真实；如果我们找不到他们，你们就知道这都是谎言。}\par

于是议会突然起身，决定选派合宜的人——敬畏天主、并知道那些人何时去世、葬在何处——去仔细查问，看看是否如他们所听见的那样。于是有十五个人前往那个地方，他们都在那些人安息时在场，埋葬时曾站在他们脚旁，并见过他们的坟墓。他们来到那里，发现坟墓是敞开的，还有许多别的坟墓也是如此，既找不到骨头的痕迹，也找不到骨灰。他们急忙返回，报告了所见之事。\par

那时，他们的全部会堂陷入极大的忧伤和困惑，彼此说：\Quote{我们该怎么办？}\Person{亚纳斯}和\Person{盖法}说：\Quote{让我们转向听说他们在的地方，派些有地位的人去，请求并恳求他们：也许他们会屈尊来见我们。}于是他们派了\Person{尼苛德摩}、\Person{若瑟}以及那三位曾看见他们的加里肋亚辣彼去，请他们屈尊来见他们。他们去了，走遍\Place{约但河}整个地区和山岭，却空手而归，没有找到他们。\par

看，忽然从\Place{阿玛肋克山}下来一大群人，大约一万二千名，是那些与主一同复活的人。虽然他们认出了许多人，但因恐惧和天使的显现而不能向他们说话；他们远远地站着，注视着并听着他们走路唱歌赞美说：\Quote{主已从死者中复活，正如祂所说的；让我们都欢欣喜乐，因为祂永远为王。}于是被派去的人惊愕不已，因恐惧而仆倒在地，并从他们得到答复：他们将在自己的家中见到\Person{卡里诺斯}和\Person{路基乌斯}。\par

他们起身前往他们的家，发现他们正在祈祷。他们进去，面伏于地，向他们致敬；站起身来，说：\Quote{天主的友人，全体犹太群众派我们来见你们，听说你们从死者中复活了，他们请求并恳求你们到他们那里去，好让我们都知道在我们这个时代所发生的伟大天主之事。}他们立刻因天主的指示起身，与他们同去，进入了他们的会堂。于是犹太群众和司祭们将法律书放在他们手中，以天主\OriginalTerm{赫罗依}{Heloi}、天主\OriginalTerm{阿多奈}{Adonai}以及法律和先知的名义起誓说：\Quote{告诉我们你们如何从死者中复活，以及在我们这个时代发生了哪些奇事，是我们从未听过在任何时代发生过的；因为恐惧之下，我们所有的骨头都已麻木干枯，大地在我们脚下震动：我们全心合力倾流了义人和圣洁的血。}\par

于是\Person{卡里诺斯}和\Person{路基乌斯}向他们示意，要他们给一张纸和墨水。他们这样做，因为圣神不允许他们向他们说话。他们就每人给了一张纸，把他们分开在各自隔开的房间里。他们用指头画了基督的十字圣号，开始在各自的纸上书写；写完后，仿佛从同一张口中，从各自的隔间里喊出：\Quote{阿们。}然后站起来，\Person{卡里诺斯}把他的纸交给\Person{亚纳斯}，\Person{路基乌斯}交给\Person{盖法}；彼此告别后，他们出去，回到自己的坟墓。\par

于是\Person{亚纳斯}和\Person{盖法}打开纸卷，各人秘密地开始读。但全体人民对此感到不满，所以大家都喊道：\Quote{请公开地向我们宣读这些文字；读完后我们将保存它们，免得这真理因故意盲目而被不洁和诡诈的人歪曲成谎言。}\Person{亚纳斯}和\Person{盖法}听了这话，颤抖起来，把纸卷交给从\Place{加里肋亚}来的\Person{辣彼阿达斯}、\Person{辣彼斐乃斯}和\Person{辣彼厄基雅斯}，他们曾宣告耶稣被提到天上。全体犹太群众信赖他们来宣读这文件。他们便读了那纸卷上的这些话：——\par

% structure: p0012 source=h2 target=subsection reason=relative-heading-rank; base=h2

\subsection{\Emph{第二章（18）}}

\Signature{我，\Person{卡里努斯}}。主耶稣基督，永生天主之子，请允许我讲述祢在阴府中所行的奇事。当我们被拘留在阴府黑暗与死影中时，忽然有巨光照耀我们，阴府和死亡的门户战栗。随后听见至高圣父之子的声音，犹如雷鸣；祂大声宣告，如此命令：\Quote{你们，诸位首领，举起门户；永远的门户，请举起；光荣的君王，基督上主，将要进来。}\par

那时，死亡的首领撒殚惊恐地逃来，对他的官吏和下面的权势说：\Quote{我的官吏，以及下面所有的权势，你们齐集，关闭门户，架上铁杠，勇敢作战，抵抗，免得他们捉住我们，用锁链囚禁我们。}于是，他所有不虔敬的官吏都惊慌失措，开始尽力关闭死亡的门户，渐渐锁上门闩和铁杠，手中握紧所有武器，发出可怕而极其恐怖的嚎叫。\par

% structure: p0015 source=h2 target=subsection reason=relative-heading-rank; base=h2

\subsection{第三章（19）}

那时撒殚对阴府说：\Quote{你准备接待我将要带下去给你的那位。}于是阴府回答撒殚说：\Quote{那声音不是别的，正是至高圣父之子的呼喊，因为大地和阴府的一切处所都因此在震撼：所以我以为我自己和我所有的牢狱现在都已敞开。但我恳求你，撒殚，一切罪恶的首领，凭你的能力和我的能力，不要把他带到我这里，免得我们想捉拿他，反被他俘获。因为单凭他的声音，我的全部权势就已如此毁灭，你想他亲自来时将会怎样？}\par

死亡的首领撒殚回答他说：\Quote{你喊叫什么？不要害怕，我古老的最邪恶的朋友，因为我已煽动犹太人反对他；我已告诉他们用掌击打他的脸，我已让他的一个门徒出卖他；他是一个非常怕死的人，因为出于恐惧他说：\InnerQuote{我的心灵忧闷得要死}；我已使他如此，以致他刚被举起钉在十字架上。}\par

于是阴府对他说：\Quote{如果他就是那位仅凭他命令的一句话，使\Person{拉匝禄}在死后四天如鹰从我怀中飞走的那位，那么他不是人性的人，而是威严的天主。我恳求你不要把他带到我这里。}撒殚对他说：\Quote{无论如何你准备自己；不要害怕，因为他已挂在十字架上，我别无他法。}于是阴府回答撒殚说：\Quote{那么，如果你别无他法，看，你的毁灭就在眼前。总之，我将被丢弃而受辱；而你，却要在我权下受折磨。}\par

% structure: p0019 source=h2 target=subsection reason=relative-heading-rank; base=h2

\subsection{第四章（20）}

天主的圣者们听到了撒殚和阴府的争吵。他们虽然彼此还未完全认识，却仍保有理智。但我们的圣祖\Person{亚当}即回答撒殚说：\Quote{死亡的首领，你为什么害怕战栗？看，主就要来了，祂现在将要摧毁你的一切诡计；你必被祂拿住，永远被捆绑。}\par

那时，所有圣者听见我们的圣祖亚当如何勇敢地逐点回答撒殚，便在喜乐中坚强起来；大家一齐跑到圣祖亚当面前，聚集在一处。于是我们的圣祖亚当凝视那整个人群，大为惊奇，不知道他们是否都是从他在世上生出来的。他拥抱所有站在他周围的人们，流下极苦的眼泪，对他的儿子\Person{舍特}说：\Quote{我儿\Person{舍特}，请向圣祖们和先知们讲述，当我生病时派你去取那怜悯之油来抹我身体时，乐园的守护者对你说了什么。}\par

于是他回答说：\Quote{当你派我到乐园门前时，我流泪祈祷恳求主，并呼求乐园的守护者把油给我。于是总领天使\Person{弥额尔}出来对我说：\Person{舍特}，你为什么哭？要知道，预先得知，你的父亲\Person{亚当}今生不会得到这怜悯之油，而是要在许多世代之后。因为天主最爱的圣子将从天降下进入世界，并由\Person{若翰}在\Place{约但河}受洗；那时你的父亲\Person{亚当}将得到这怜悯之油，以及所有相信祂的人。凡相信祂的人，他们的王国将永存。}\par

% structure: p0023 source=h2 target=subsection reason=relative-heading-rank; base=h2

\subsection{第五章（21）}

那时，所有圣者又听到这话，欢欣鼓舞。有一位站在周围的人，名叫\Person{依撒意亚}，大声呼喊，如雷吼道：\Quote{圣祖亚当，以及所有在场的人，请听我的宣告。当我在世时，藉圣神的教导，在预言中我歌颂了这光：\InnerQuote{坐在黑暗中的百姓，看见了一道大光；住在死影之地的人，光已升起。}}听到这话，圣祖亚当和众人转身问他：\Quote{你是谁？因为你说的是真的。}他又接着说：\Quote{我的名字是\Person{依撒意亚}。}\par

接着又有一位出现在他旁边，仿佛一位隐修士。他们问他说：\Quote{你是谁？你的身体为什么有这样的相貌？}他坚定地回答说：\Quote{我是洗者\Person{若翰}，至高者的声音和先知。我走在这位主的前面，为使荒芜崎岖之地成为平坦道路。我用手指指明并将主羔羊、天主子显示给\Place{耶路撒冷}的居民。我在\Place{约但河}为他施洗。我听见从天上来的圣父的声音，如雷响在他上面，宣告：\InnerQuote{这是我的爱子，我所喜悦的。}我从他得到答复，他要下到阴府。}\par

圣祖亚当听见这话，大声呼喊说：\OriginalTerm{阿肋路亚}{Alleluia}！意思是：\Quote{主确实要来了。}\par

% structure: p0027 source=h2 target=subsection reason=relative-heading-rank; base=h2

\subsection{第六章（22）}

此后，另有一人站在那里，仿佛带着某种帝王的标记，名叫达味，极为卓越，如此呼喊说：当我在世时，我曾向百姓启示天主的仁慈和祂的眷顾，预言未来的喜乐，历世历代说：\BibleQuote{愿他们称谢上主的仁慈和祂向人子所行的奇事，因祂打碎了铜门，砍断了铁闩。}于是圣祖和先知们开始彼此相认，各人引用自己的预言。\par

随后，圣耶肋米亚查考自己的预言，对圣祖和先知们说：当我在世时，我曾预言天主之子，说祂在地上显现，并与世人同住。\par

于是，众圣徒在主的光明中欢欣，目睹父亚当，并听到所有圣祖和先知的回应，呼喊说：\BibleQuote{阿肋路亚！奉主名而来的当受赞美！}以致\OriginalTerm{撒殚}{Satan}战栗，寻找逃脱之路。但他不能，因为\OriginalTerm{阴府}{Hades}及其属下将他捆绑在阴间，并严密看守。他们对他说：你为什么战栗？我们绝不允许你从这里出去。但要从你每日攻击的那位手中接受这应得的惩罚；否则，你要知道，你被他捆绑，将受我约束。\par

% structure: p0031 source=h2 target=subsection reason=relative-heading-rank; base=h2

\subsection{第7章（23）}

再次传来至高圣父之子的声音，如同大雷之声，说：\BibleQuote{众门哪，抬起头来！永久的门户啊，你们要被举起！光荣的君王要进来！}于是\OriginalTerm{撒殚}{Satan}和\OriginalTerm{阴府}{Hades}呼喊说：\BibleQuote{谁是这位光荣的君王？}他们以主的声音回答说：\BibleQuote{是上主，大能而有力的上主，是战场上大能的上主。}\par

这声音之后，来了一个人，形貌如强盗，肩上背着十字架，在门外呼喊说：给我开门，让我进去！\OriginalTerm{撒殚}{Satan}为他稍微开了门，将他带进自己的居所，随后又把门关上。众圣徒清清楚楚地看见他，立刻对他说：你的形貌像个强盗。告诉我们你背上扛的是什么。他谦卑地回答说：我确实是个十足的强盗；犹太人把我挂在十字架上，与我的主耶稣基督，至高圣父之子一同悬在木架上。总之，我是来为祂传报的；祂确实紧随我后而来。\par

这时，圣达味对\OriginalTerm{撒殚}{Satan}怒火中烧，大声呼喊说：\BibleQuote{打开你的门，最卑鄙的恶徒，让光荣的君王进来！}同样，天主所有的圣徒都起来攻击\OriginalTerm{撒殚}{Satan}，要抓住他，将他瓜分。又有声音在内部响起：\BibleQuote{众门哪，抬起头来！永久的门户啊，你们要被举起！光荣的君王要进来！}\OriginalTerm{阴府}{Hades}和\OriginalTerm{撒殚}{Satan}在那清楚的声音中再次问道：\BibleQuote{这位光荣的君王是谁？}那奇妙的声音对他们说：\BibleQuote{万军的上主，祂是光荣的君王。}\par

% structure: p0035 source=h2 target=subsection reason=relative-heading-rank; base=h2

\subsection{第8章（24）}

看，忽然\OriginalTerm{阴府}{Hades}战栗，死亡的门户和门闩碎裂，铁闩折断落地，一切豁然开朗。\OriginalTerm{撒殚}{Satan}站在中间，惊惶沮丧，双脚被锁链捆绑。看，主耶稣基督从高天带着光明的照耀而来，慈悲、伟大而谦卑，手中拿着锁链，捆绑\OriginalTerm{撒殚}{Satan}的颈项；又将他的双手反绑，将他背朝下摔进\OriginalTerm{塔尔塔罗斯}{Tartarus}，将圣足踏在他的喉咙上说：你世世代代行了诸多恶事，从未安息。今天我将你交付永火。又突然召唤\OriginalTerm{阴府}{Hades}，命令他说：将这最邪恶、最不虔敬者拿去，看守在我吩咐你的那天之前。他一接过来，便同主足下被投入深渊之底。\par

% structure: p0037 source=h2 target=subsection reason=relative-heading-rank; base=h2

\subsection{第9章（25）}

然后，主耶稣，万民的救主，慈悲而极其温良，慈祥地招呼亚当，对他说：愿平安归于你，亚当，与你的众子，直到无穷的世世代代！阿们。于是父亚当俯伏在主脚下，被扶起后，亲吻了祂的双手，流了许多眼泪，向众人作证说：看，这双手塑造了我！又对主说：光荣的君王，你来了，拯救世人，领他们进入你永生的国度。同样，我们的母亲厄娃也俯伏在主脚下，被扶起，亲吻了祂的双手，流了许多眼泪，向众人作证说：看，这双手造了我！\par

于是，众圣徒敬拜祂，呼喊说：\BibleQuote{奉主名而来的当受赞美！主上帝光照了我们——阿们——直到永远。阿肋路亚，永永远远！赞美、尊荣、权能、光荣！因为你从高处降临眷顾了我们。}他们不断歌唱阿肋路亚，因祂的荣耀同乐，一同拥到主的手下。然后救主仔细查问一切，抓住\OriginalTerm{阴府}{Hades}，立即将一些投入\OriginalTerm{塔尔塔罗斯}{Tartarus}，又带领一些与祂一同升到上界。\par

% structure: p0040 source=h2 target=subsection reason=relative-heading-rank; base=h2

\subsection{第10章（26）}

于是，天主所有的圣徒恳求主，将祂圣十字架的记号留在阴间作为胜利的标志，使阴间最不虔敬的官吏不得拘留任何主所赦免的人为罪犯。事情就这样成就了。主将祂的十字架立在\OriginalTerm{阴府}{Hades}中间，作为胜利的记号，它将留存直到永远。\par

然后我们所有人与主一同从那里出来，将\OriginalTerm{撒殚}{Satan}和\OriginalTerm{阴府}{Hades}留在\OriginalTerm{塔尔塔罗斯}{Tartarus}中。我们和许多其他人奉命要肉身复活，在世上为主耶稣基督的复活以及阴间所发生的事作证。\par

最亲爱的弟兄们，这些就是我们所见的事，并且因受你们恳求而作证，那位为我们死而复活的主亲自作证；因为经上如何记载，一切都如何应验了。\par

% structure: p0044 source=h2 target=subsection reason=relative-heading-rank; base=h2

\subsection{第11章（27）}

当那文书读完并宣读完毕时，所有听见的人都俯伏在地，痛哭流涕，并狠狠地捶胸，呼喊着，齐声说：\Quote{我们有祸了！为何这事临到我们这些可怜的人身上？\Person{比拉多}逃跑了，\Person{亚纳斯}和\Person{盖法}逃跑了，司铎和肋未人也逃跑了；连犹太百姓也哭着说：\InnerQuote{我们有祸了！我们流了圣血在地上。}}\par

因此，三天三夜，他们完全没有尝过饼和水，也没有任何人回\Place{会堂}去。但第三天，公议会再次聚集，\Person{路济乌斯}的另一份文书也被宣读，发现与\Person{卡里努斯}所写的文书毫无出入，连一个字母也不差。于是\Place{会堂}困惑了；他们全体哀哭了四十昼夜，等候来自\OriginalTerm{天主}{God}的毁灭和\OriginalTerm{天主}{God}的报复。但祂是充满深情、最崇高的怜悯者，没有立即毁灭他们，而是慷慨地赐予他们悔改的机会。但他们却不配归向\OriginalTerm{主}{Lord}。\par

这些是\Person{卡里努斯}和\Person{路济乌斯}的见证，最亲爱的弟兄们，关于\Person{基督}\OriginalTerm{天主子}{Son of God}，以及祂在阴府中的神圣事迹；愿我们众人将赞颂和荣耀归于祂，直到无穷无尽永世。阿们。\par

\SourceRecord{Translated by Alexander Walker.；Ante-Nicene Fathers；Vol. 8.；Edited by Alexander Roberts, James Donaldson, and A. Cleveland Coxe.；Buffalo, NY: Christian Literature Publishing Co.,；1886.；<http://www.newadvent.org/fathers/08072c.htm>.}
